\begin{textsample}{NEG Dim 1 – 2021\_c6 – Score -4.00 – 2021\_c6/enade\_2021\_mus\_8d85dac343.txt}  \label{ex:f1_neg_235}
CAPÍTULO 6 ANÁLISE TÉCNICA DA PROVA Este capítulo tem por objetivo apresentar o desempenho dos estudantes concluintes de Música ( Licenciatura ) no Enade/2021.
Para isso, foram calculadas as estatísticas básicas da prova em seu todo ( seção 6.1.1 ), bem como as estatísticas relacionadas aos Componentes de Formação Geral ( seção 6.1.2 ) e de Conhecimento Específico da Área ( seção 6.1.3 ).
Dadas as suas características, foram analisadas, em separado, as questões objetivas ( seção 6.2 ) e as questões discursivas ( seção 6.3 ).
Tomando -se como base as duas Questões Discursivas do Componente de Formação Geral, nas seções 6.3.1.6 a 6.3.1.8, são apresentados os resultados e os comentários sobre a correção das respostas em relação à Língua Portuguesa.
Para as questões objetivas, foram disponibilizados os índices de Facilidade e de Discriminação Ponto-Bisserial, também em separado, para os Componentes de Formação Geral ( seção 6.2.1 ) e de Conhecimento Específico ( seção 6.2.2 ).
De \textbf{cada} componente, uma das questões foi escolhida para exemplificar a análise gráfica, relacionando -se as alternativas escolhidas pelos estudantes ( inclusive o gabarito ) com o número de acertos no componente.
No Anexo I, apresenta -se a íntegra da análise gráfica para todas as questões objetivas.
Para \textbf{cada} uma das questões discursivas, os conteúdos dos tipos mais comuns de respostas dos estudantes são apresentados e comparados com o padrão de resposta esperado ( ver Anexo IX com o padrão de respostas ).
Nas tabelas, constam as seguintes estatísticas das notas¹⁹ : média do desempenho na prova, erro padrão da média, desvio padrão, nota mínima, mediana e nota máxima para \textbf{cada} um de seus componentes.
Tais estatísticas contemplam o total de estudantes concluintes da Área de Música ( Licenciatura ) inscritos e presentes à prova do Enade/2021, tendo em vista agregações, ou por Grandes Regiões e o país como um todo, ou por Categoria Administrativa, Organização Acadêmica da IES e Modalidade de Ensino.
Em relação aos gráficos de distribuição de notas, o intervalo considerado foi de 10 unidades, aberto à esquerda e fechado à direita, com exceção do primeiro intervalo, [ 0 ; 10 ], fechado em ambos os extremos.
Para os gráficos de distribuição das notas das questões discursivas, foram consideradas mais duas categorias : questão em branco²⁰ e nota zero. ¹⁹ Uma definição dessas estatísticas pode ser encontrada no Glossário. ²⁰ Nesse grupo \textbf{estão} incluídas também as respostas classificadas como nulas ou desconsideradas. 6.1 ESTATÍSTICAS BÁSICAS DA PROVA Esta seção apresenta estatísticas selecionadas e histogramas da nota geral ( 6.1.1 ) e de \textbf{cada} componente : Formação Geral ( 6.1.2 ) e Conhecimento Específico ( 6.1.3 ).
São também apresentadas estatísticas selecionadas de subpopulações, caracterizadas por Grande Região, Categoria Administrativa, Organização Acadêmica e Modalidade de Ensino. 6.1.1 ESTATÍSTICAS BÁSICAS GERAIS Na Tabela 6.1, são apresentadas as Estatísticas Básicas da prova, por Grande Região, dos estudantes concluintes de Música ( Licenciatura ).
A Média das notas da prova, como um todo ( nas seções seguintes serão analisados os Componentes de Formação Geral e de Conhecimento Específico ), foi 42,3, sendo que os estudantes da região Norte obtiveram a Média mais baixa ( 36,9 ), e os da região Centro-Oeste obtiveram a Média mais alta ( 45,6 ).
As demais Médias foram : 38,9, na região Nordeste, 43,7, na região Sudeste e 44,6, na região Sul.
O Desvio padrão para o Brasil, como um todo, foi 13,8, o mesmo da região Sudeste.
O menor Desvio padrão foi encontrado na região Norte ( 12,5 ), indicando uma dispersão um pouco menor das notas dessa última região.
A região que obteve a maior nota Máxima foi a Sudeste ( 83,2 ), ao passo que a região que atingiu a menor nota Máxima foi a Centro-Oeste ( 72,2 ).
A Mediana do Brasil, como um todo, foi 41,8, sendo a maior Mediana obtida na região Centro-Oeste ( 47,2 ), e a menor obtida na região Norte ( 35,8 ).
A nota Mínima do Brasil foi obtida na região Nordeste, onde foi zero.
A maior nota Mínima, entre as cinco Grandes Regiões, foi 13,2, na região Norte, as demais regiões obtiveram as seguintes notas Mínima : 8,5 ( Sudeste ), 2,3 ( Sul ) e 12,7 ( Centro-Oeste ).
Considerando -se as notas segundo Grande Região, observa -se que existe diferença estatisticamente significativa ao nível de 95 \% entre a Média obtida pela região Norte em relação às Médias das regiões Sudeste, Sul e Centro-Oeste, e entre as Médias obtidas pela região Nordeste em relação às obtidas por essas mesmas três regiões.
O comportamento das notas dos estudantes de todo o Brasil pode ser observado no Gráfico 6.1, que apresenta um histograma com a distribuição das mesmas : uma distribuição unimodal com a moda no intervalo ( 30 ; 40 ].
Na Tabela 6.2, são apresentadas informações referentes à Média da nota final, desagregadas por Categoria Administrativa, por Organização Acadêmica e por Modalidade de Ensino.
A Média das notas da prova, como um todo, foi 42,3.
Em relação a Categoria Administrativa, os estudantes das IES Públicas obtiveram Média mais alta ( 43,0 ) que a Média nacional, e os das IES Privadas obtiveram Média mais baixa ( 41,4 ).
Observa -se que existe diferença estatisticamente significativa entre as Médias das notas das IES Públicas e as das Privadas.
A diferença entre as Médias das regiões Centro-Oeste e Norte ( 8,7 ), a maior e a menor Média, é superior à diferença entre as Médias das IES Públicas e Privadas ( 1,6 ), caracterizando uma maior diversidade regional do que administrativa.
O Desvio padrão para as IES Públicas ( 14,3 ) foi superior ao do Brasil, como um todo ( 13,8 ), indicando uma dispersão um pouco maior das notas nesta Categoria Administrativa.
No tocante à Organização Acadêmica, as Universidades ( 43,3 ) obtiveram Média um ponto mais alta que a nacional ( 42,3 ).
A Média dos Centros Universitários ( 40,9 ), das Faculdades ( 40,2 ) e dos CEFET/IF ( 41,5 ) foram menores do que a nacional.
Constata -se que existe diferença estatisticamente significativa, ao nível de 95 \%, entre as Médias das Universidades e as Médias dos Centros Universitários e das Faculdades.
Faculdades ( 40,2 ) e dos CEFET/IF ( 41,5 ) foram menores do que a nacional.
Constata -se que existe diferença estatisticamente significativa, ao nível de 95 \%, entre as Médias das Universidades e as Médias dos Centros Universitários e das Faculdades.
A Média da Modalidade Educação a Distância ( 39,0 ) foi inferior à Média da Educação Presencial ( 44,6 ), e a diferença entre elas é estatisticamente significativa, ao nível de 95 \%. 6.1.2 ESTATÍSTICAS BÁSICAS NO COMPONENTE DE FORMAÇÃO GERAL Na Tabela 6.3, são apresentadas as Estatísticas Básicas em relação ao componente da prova que avalia a Formação Geral dos estudantes concluintes.
A nota de Formação Geral é obtida levando -se em conta os acertos nas oito questões objetivas e a nota média obtida nas Questões Discursivas 1 e 2, sendo a nota de \textbf{cada} questão discursiva obtida levando -se em conta 80 \% da nota relativa ao conteúdo e 20 \% da nota relativa ao desempenho linguístico.
Os estudantes de todo o Brasil obtiveram desempenho médio de 35,7.
Quanto à variabilidade, o Desvio padrão das notas dos estudantes do Brasil foi 15,4.
A maior Média foi obtida na região Sul ( 38,1 ) e a menor, na região Norte ( 32,2 ).
As demais Médias foram : 33,1, na região Nordeste, 36,6, na região Sudeste e 36,0, na região Centro-Oeste.
Já o maior Desvio padrão foi obtido na região Norte ( 16,4 ), e o menor, na região Sul ( 14,7 ).
Os demais desvios foram : 14,9, na região Nordeste, 15,6, na região Sudeste e 15,3, na região Centro-Oeste.
A maior nota no Componente de Formação Geral da prova do Enade/2021 foi 84,4, obtida por, pelo menos, um estudante na região Norte.
A menor nota Máxima foi obtida na região Centro-Oeste ( 72,3 ).
A Mediana do Brasil, como um todo, foi 35,3, sendo a menor Mediana encontrada na região Norte ( 30,0 ), e a maior, nas regiões Sul e Centro-Oeste ( 37,5 ).
A nota Mínima nessa parte foi zero em todas as regiões.
Considerando -se as notas, segundo Grande Região, observa -se que existe diferença estatisticamente significativa apenas entre as Médias das notas do Componente de Formação Geral das regiões Sudeste e Sul em relação às Médias das regiões Norte e Nordeste.
No Gráfico 6.2, é possível verificar a avaliação do desempenho dos estudantes do Componente de Formação Geral, a partir do histograma da distribuição das notas correspondentes.
A distribuição é unimodal, com moda no intervalo ( 20 ; 30 ] e frequência de 22,8 \%.
Esse intervalo é imediatamente anterior ao modal da distribuição de notas da prova, como um todo ( Gráfico 6.1 ).
O intervalo modal é seguido de perto pelo intervalo ( 30 ; 40 ], com frequência de 22,5 \%.
Nota -se, ainda, que, no Gráfico 6.2, as notas apresentam dispersão um pouco maior do que as no Gráfico 6.1 ( distribuição das notas da prova ), confirmada pela comparação dos desvios padrões : 13,8 para a nota da prova, como um todo, e 15,4 para o Componente de Formação Geral.
Na Tabela 6.4, são apresentadas as informações referentes ao desempenho dos concluintes do Componente de Formação Geral, em diferentes agregações : Categoria Administrativa, Organização Acadêmica e Modalidade de Ensino.
Observa -se que não existe diferença estatisticamente significativa entre as Médias dos tipos de Categoria Administrativa, que têm valores muito próximos : a maior Média obtida por estudantes de IES Privadas ( 35,7 ) e a menor, pelos de IES Públicas ( 35,6 ).
Nota -se, no entanto, que há diferença estatisticamente significativa entre a Média das Universidades ( 36,2, a maior ) em relação à Média das Faculdades ( 32,3, a menor ) e à Média dos CEFET/IF ( 32,6 ).
A Média dos Centros Universitários foi 35,4.
No que se refere a Modalidade de Ensino, a Educação Presencial apresentou maior Média do que a Educação a Distância ( 36,4 e 34,6, respectivamente ).
Tais Médias apresentam diferença estatisticamente significativa. 6.1.3 ESTATÍSTICAS BÁSICAS DO COMPONENTE DE CONHECIMENTO ESPECÍFICO Na Tabela 6.5, são apresentadas as Estatísticas Básicas referentes ao Componente de Conhecimento Específico da Área de Música ( Licenciatura ).
A nota do componente de Conhecimentos Específicos leva em conta as notas da parte objetiva e a média das notas obtidas nas três questões discursivas parte discursiva da Área.
A Média do desempenho dos estudantes do Brasil, como um todo, foi 44,5.
A maior Média foi obtida na região Centro-Oeste ( 48,8 ), e a menor, na região Norte ( 38,4 ).
As demais Médias foram : 40,9 na região Nordeste, 46,0 na região Sudeste e 46,7 na região Sul.
Quanto à variabilidade das notas, o Desvio padrão do Brasil, como um todo, foi 15,1, sendo o maior Desvio padrão observado na região Sul ( 15,1 ) e o menor, na região Norte ( 12,8 ).
Os demais desvios foram : 14,9 na região Nordeste, 15,0 na região Sudeste e 14,2 na região Centro-Oeste.
A Mediana das notas dos estudantes de todo o Brasil foi 43,5.
A maior Mediana ocorreu na região Centro-Oeste ( 51,1 ), e a menor, na região Norte ( 38,1 ).
As demais Medianas foram : 40,2 na região Nordeste, 45,1 na região Sudeste e 45,7 na região Sul.
A nota Máxima do Brasil, como um todo, foi 88,1, sendo obtida por pelo menos um estudante na região Sudeste.
As demais notas Máximas foram : 81,3 na região Norte, 85,4 na região Nordeste, 87,4 na região Sul e 76,0 na região Centro-Oeste.
A nota Mínima foi zero na região Nordeste.
As demais notas Mínimas foram : 12,6 na região Norte, 4,2 na região Sudeste, 3,1 na região Sul e 15,7 na região Centro-Oeste.
Observa -se que não existe diferença estatisticamente significativa entre as regiões Sudeste, Sul e Centro-Oeste na análise das Médias das notas do Componente de Conhecimento Específico, mas existe diferença estatisticamente significativa entre \textbf{cada} uma dessas três regiões e as regiões Norte e Nordeste.
Assim como nos Gráficos 6.1 e 6.2, no Gráfico 6.3, apresentado a seguir, é possível ser feita uma avaliação do desempenho de concluintes de Música ( Licenciatura ), em relação ao Componente de Conhecimento Específico, com um histograma da distribuição das notas correspondentes.
Esta também é uma distribuição unimodal, e o grupo modal é o ( 40 ; 50 ], um intervalo acima do grupo modal da prova, como um todo, e dois intervalos acima do grupo modal para a Formação Geral.
A frequência do intervalo modal, 22,9 \%, é bem próxima a do intervalo ( 30 ; 40 ], com frequência 22,7 \%.
Na Tabela 6.6, é apresentada uma comparação dos resultados em relação a Categoria Administrativa, a Organização Acadêmica e a Modalidade de Ensino, agora levando -se em conta o desempenho de estudantes do Componente de Conhecimento Específico da prova.
No que se refere a Organização Acadêmica, a maior Média foi a das Universidades ( 45,7 ), vindo a seguir a dos CEFET/IF ( 44,5 ) e, por fim, a dos Centros Universitários e a das Faculdades ( 42,8, ambos ).
O maior Desvio padrão foi o das Universidades ( 15,7 ), maior que o valor para o Brasil, como um todo.
As Universidades obtiveram a maior nota Máxima ( 88,1 ).
Os Centros Universitários obtiveram nota Máxima 84,1, as Faculdades, 79,9, e os CEFET/IF, 85,4.
As Medianas foram : 45,1 nas Universidades, 42,1 nos Centros Universitários, 41,2 nas Faculdades e 42,2 nos CEFET/IF.
A nota Mínima foi zero para as Universidades, 4,2 para os Centros Universitários, 11,4 para as Faculdades, e 18,9 para os CEFET/IF.
Observa -se que existe diferença estatisticamente significativa, ao nível de 95 \%, no Componente de Conhecimento Específico, apenas entre a Média das Universidades e a dos Centros Universitários.
Quanto a Categoria Administrativa, observa -se um comportamento semelhante àquele da prova, como um todo, ou seja, existe diferença estatisticamente significativa entre as Médias das IES Públicas ( 45,4 ) e as das IES Privadas ( 43,3 ).
Neste caso, também, a maior Média foi obtida por estudantes de IES Públicas de ensino.
Quanto ao comportamento das Médias, no que toca a Modalidade de Ensino, há diferença estatisticamente significativa entre a Média da Educação Presencial ( 47,3 ) e a da Educação a Distância ( 40,5 ). 6.2 ANÁLISE DAS QUESTÕES OBJETIVAS Esta seção apresenta estatísticas selecionadas e histogramas das Questões Objetivas de Formação Geral ( 6.2.1 ) e de Conhecimento Específico ( 6.2.2 ).
São, também, apresentadas e comparadas as médias das subpopulações caracterizadas por Grande Região. 21 6.2.1 COMPONENTE DE FORMAÇÃO GERAL Médias das IES Públicas ( 45,4 ) e as das IES Privadas ( 43,3 ).
Neste caso, também, a maior Média foi obtida por estudantes de IES Públicas de ensino.
Quanto ao comportamento das Médias, no que toca a Modalidade de Ensino, há diferença estatisticamente significativa entre a Média da Educação Presencial ( 47,3 ) e a da Educação a Distância ( 40,5 ). 6.2 ANÁLISE DAS QUESTÕES OBJETIVAS Esta seção apresenta estatísticas selecionadas e histogramas das Questões Objetivas de Formação Geral ( 6.2.1 ) e de Conhecimento Específico ( 6.2.2 ).
São, também, apresentadas e comparadas as médias das subpopulações caracterizadas por Grande Região. 21 6.2.1 COMPONENTE DE FORMAÇÃO GERAL Na Tabela 6.7, são apresentadas as Estatísticas Básicas relativas às oito Questões Objetivas do Componente da prova que abrange a Formação Geral dos estudantes.
A Média do Brasil foi 41,1.
A menor Média foi encontrada na região Norte ( 37,7 ), e a maior, na região Centro-Oeste ( 43,4 ).
As demais Médias foram : 39,4 na região Nordeste, 41,6 na região Sudeste e 42,8 na região Sul.
O Desvio padrão do Brasil foi 19,9, sendo o maior Desvio padrão encontrado nas regiões Norte e Centro-Oeste ( 21,0 ), e o menor Desvio padrão encontrado na região Sul ( 18,7 ).
Nas demais regiões, os valores do Desvio padrão foram 19,8 ( Nordeste ) e 20,1 ( Sudeste ).
As Medianas do Brasil, como um todo, e de quatro regiões ( Norte, Nordeste, Sudeste e Sul ) foram iguais a 37,5.
Na região Centro-Oeste a Mediana foi 50,0.
A nota Máxima ( 100,0 ) foi alcançada nas regiões Norte e Sudeste, e, nas demais, a nota Máxima foi 87,5.
As notas Mínimas ( 0,0 ) foram iguais para todas as regiões.
Na Tabela 6.8, são apresentados o Índice de Facilidade e o Índice de Discriminação ( Ponto-Bisserial ) para \textbf{cada} uma das Questões Objetivas do Componente de Formação Geral.
Quanto ao Índice de Facilidade, foram usadas as seguintes cores para diferenciar o nível de dificuldade da questão :  Azul para as questões classificadas com índice Muito fácil ( \&gt ; =0,86 ), verde para as questões classificadas com índice Fácil ( 0,61 a 0,85 ), amarelo para as questões classificadas com Médio ( 0,41 a 0,60 ), vermelho para as questões classificadas com Difícil ( 0,16 a 0,40 ) e roxo para as questões classificadas com Muito difícil ( \&lt ; =0,15 ).
Já quanto ao Índice de Discriminação, foram usadas as seguintes cores para qualificar a questão :  As questões classificadas com índice Fraco receberam a cor vermelha ( \&lt ; =0,19 ), as classificadas com Médio receberam a cor amarela ( 0,20 a 0,29 ), as classificadas com Bom receberam a cor verde ( 0,30 a 0,39 ) e as classificadas com Muito bom ( \&gt ; =0,40 ) receberam a cor azul.
As Questões Objetivas do Componente de Formação Geral, segundo o Índice de Facilidade, foram assim avaliadas : nenhuma teve o Índice de Facilidade classificado como Muito fácil, e uma questão foi tida como Fácil, por ter índice de acertos 0,61.
Quatro questões foram consideradas com índice de dificuldade Médio, situando -se no intervalo entre 0,41 e 0,60 do Índice de Facilidade, ou seja, houve entre 41 \% e 60 \% de acertos, enquanto três questões foram classificadas na categoria Difícil, apresentando percentual de acertos entre 16 \% e 40 \%.
Nenhuma questão apresentou percentual de acertos igual ou inferior a 15 \%, razão pela qual não houve questão classificada como Muito Difícil. 21 As fórmulas utilizadas para o coeficiente de correlação ponto-bisserial e para o índice de facilidade \textbf{estão} apresentadas no capítulo 1. foi alcançada nas regiões Norte e Sudeste, e, nas demais, a nota Máxima foi 87,5.
As notas Mínimas ( 0,0 ) foram iguais para todas as regiões.
Na Tabela 6.8, são apresentados o Índice de Facilidade e o Índice de Discriminação ( Ponto-Bisserial ) para \textbf{cada} uma das Questões Objetivas do Componente de Formação Geral.
Quanto ao Índice de Facilidade, foram usadas as seguintes cores para diferenciar o nível de dificuldade da questão : • Azul para as questões classificadas com índice Muito fácil ( \&gt ; =0,86 ), verde para as questões classificadas com índice Fácil ( 0,61 a 0,85 ), amarelo para as questões classificadas com Médio ( 0,41 a 0,60 ), vermelho para as questões classificadas com Difícil ( 0,16 a 0,40 ) e roxo para as questões classificadas com Muito difícil ( \&lt ; =0,15 ).
Já quanto ao Índice de Discriminação, foram usadas as seguintes cores para qualificar a questão : • As questões classificadas com índice Fraco receberam a cor vermelha ( \&lt ; =0,19 ), as classificadas com Médio receberam a cor amarela ( 0,20 a 0,29 ), as classificadas com Bom receberam a cor verde ( 0,30 a 0,39 ) e as classificadas com Muito bom ( \&gt ; =0,40 ) receberam a cor azul.
As Questões Objetivas do Componente de Formação Geral, segundo o Índice de Facilidade, foram assim avaliadas : nenhuma teve o Índice de Facilidade classificado como Muito fácil, e uma questão foi tida como Fácil, por ter índice de acertos 0,61.
Quatro questões foram consideradas com índice de dificuldade Médio, situando -se no intervalo entre 0,41 e 0,60 do Índice de Facilidade, ou seja, houve entre 41 \% e 60 \% de acertos, enquanto três questões foram classificadas na categoria Difícil, apresentando percentual de acertos entre 16 \% e 40 \%.
Nenhuma questão apresentou percentual de acertos igual ou inferior a 15 \%, razão pela qual não houve questão classificada como Muito Difícil.
O Índice de Facilidade variou de 0,18 a 0,61, e o de Discriminação, de 0,26 a 0,49.
Seis questões, quanto ao Índice de Discriminação, foram tidas como Muito bom, uma questão foi tida como Bom, e uma como Médio.
Na Tabela 6.9, é apresentada a distribuição das questões, levando -se em conta, simultaneamente, a classificação dos dois índices.
As seis questões que tiveram Índice de Discriminação Muito bom figuraram em três níveis de dificuldade Fácil, Médio e Difícil : uma classificada na categoria Fácil ( a Questão 7 ) do Índice de Facilidade, uma classificada na categoria Difícil ( a Questão 4 ) do Índice de Facilidade e quatro, na categoria Médio ( Questões 1, 3, 5 e 6 ) do Índice de Facilidade.
Em particular, a Questão 1 foi a que apresentou o maior poder discriminatório, com índice 0,49, e foi considerada Médio em termos de Facilidade, com uma proporção de 0,46 acertos.
O máximo de acertos foi alcançado pela Questão 7, com um Índice de Facilidade de 0,61 e Índice de Discriminação 0,48.
A Questão 4 apresentou Índice de Facilidade 0,29 e Índice de Discriminação 0,43.
As Questões 3, 5 e 6 apresentaram, respectivamente, Índices de Facilidade 0,50, 0,47 e 0,55, e Índices de Discriminação 0,48, 0,44 e 0,43.
A Questão 2 foi classificada como Difícil ( 21 \% de acertos ) e com boa discriminação ( 0,36 ) e a Questão 8, Difícil ( 18 \% de acertos ) e com índice de discriminação Médio ( 0,26 ).
No Gráfico 6.4, para exemplificar, analisa -se o comportamento da Questão de número 3 de Formação Geral.
Trata -se de uma questão considerada com nível Médio de dificuldade, em relação à Facilidade, e a que obteve o segundo maior Índice de Discriminação 0,48.
Neste gráfico, \textbf{cada} uma das cinco curvas representa o percentual de respostas em determinada alternativa da questão, em função do número de acertos dos estudantes nessa parte da prova ( Formação Geral/Múltipla Escolha ), antes de possíveis eliminações pelo critério Ponto-Bisserial.
A curva em verde corresponde à alternativa B, a alternativa correta para essa questão.
Observa -se que entre os estudantes com menor número de acertos, nessa parte do exame, a situação mais frequente foi a escolha da alternativa E ( em vermelho ), incorreta.
Por exemplo, entre os estudantes que acertaram duas respostas, 45,7 \% escolheram a alternativa E ( em vermelho ), 29,0 \% escolheram a alternativa B ( em verde, gabarito ), 9,5 \% escolheram a alternativa D ( em laranja ), 9,4 \% escolheram a alternativa A ( em azul ) e 6,1 \% escolheram a alternativa C ( em preto ).
Dentre os que acertaram duas respostas entre as questões de múltipla escolha de Formação Geral, nenhum deixou a questão em branco e 0,3 \% marcou mais de uma alternativa, invalidando a questão.
À medida que o número de acertos aumenta, indicando desempenho melhor nessa parte da prova, aumenta, concomitantemente, a proporção de estudantes que selecionaram a alternativa correta B, atingindo 100 \% para os estudantes com oito acertos.
Essa análise permite verificar como a questão discriminou os grupos de desempenho, justificando -se o alto índice ( 0,48 ) obtido na questão.
Cumpre notar que não é possível inferir deste gráfico nem o Índice de Facilidade, que seria uma média da proporção ponderada pela quantidade de estudantes com \textbf{cada} uma das notas, nem o Índice de Discriminação Ponto-Bisserial, por razão equivalente.
No caso extremo, no qual a grande concentração dos acertos dos estudantes fosse abaixo de quatro, o Índice de Facilidade seria obrigatoriamente um pouco acima de 60,0 \% ( neste exemplo ).
Caso a concentração fosse em seis acertos ou mais, o índice seria obrigatoriamente acima de 80,0 \%.
Os gráficos relativos às demais questões de Formação Geral constam do Anexo I. 6.2.2 COMPONENTE DE CONHECIMENTO ESPECÍFICO Na Tabela 6.10, são apresentadas as Estatísticas Básicas em relação às Questões Objetivas do Componente de Conhecimento Específico da prova de Música ( Licenciatura ) por Grande Região.
As notas foram calculadas com base em 27 Questões Objetivas de Conhecimentos Específicos.
A Média do Brasil deste componente foi 43,9.
A menor Média foi observada na região Norte ( 38,2 ), e a maior, na região Centro-Oeste ( 48,9 ).
O Desvio padrão de todo o Brasil foi 15,6, sendo o menor Desvio padrão encontrado na região Norte ( 13,7 ), e o maior, na região Sul ( 15,7 ).
A Mediana de todo o Brasil foi 44,4, o mesmo valor da Mediana encontrada nas regiões Sudeste e Sul.
Nas regiões Norte ( 37,0 ) e Nordeste ( 40,7 ), o valor da Mediana foi menor, e, na região Centro-Oeste ( 48,1 ), foi maior.
A nota Máxima da prova foi 88,9, obtida nas Questões Objetivas do Componente de Conhecimento Específico, por, pelo menos, um estudante em três regiões Norte, Sudeste e Sul.
Na região Nordeste, a nota Máxima foi 85,2 e, na região Centro-Oeste, 77,8, a menor nota Máxima.
A nota Mínima foi zero, obtida por pelo menos um estudante na região Nordeste nas Questões Objetivas do Componente de Conhecimento Específico, e nas demais regiões, foi 11,1 ( Norte ), 3,7 ( Sudeste e Sul ) e 18,5 ( Centro-Oeste ).
Na Tabela 6.11, são apresentados os Índices de Facilidade e de Discriminação ( Ponto-Bisserial ) das Questões Objetivas do Componente de Conhecimento Específico, para os estudantes de Música ( Licenciatura ).
Para facilitar a diferenciação das questões, serão usadas as mesmas cores da Tabela 6.8 para as diferentes classificações dos Índices de Facilidade e de Discriminação.
As análises a seguir dizem respeito às 27 questões.
A partir do Índice de Facilidade obtido, observa -se que, das 27 questões válidas, uma questão objetiva da prova foi considerada Muito Difícil e 11 foram classificadas como Difícil.
Três questões foram classificadas como Fácil, e outras 11 consideradas como Médio.
Uma questão foi classificada como Muito fácil.
Já quanto aos Índices de Discriminação das Questões Objetivas do Componente de Conhecimento Específico da prova, obtém -se como resultado a seguinte classificação : dez das 27 questões válidas foram consideradas boas, enquanto seis delas tiveram Índice de Discriminação Muito bom.
Assim, para 16 em 27 questões, os Índices de Discriminação foram Bom ou Muito bom.
Dentre as demais, 11 delas foram classificadas como Médio, e nenhuma como Fraco.
Constata -se, assim, que a prova – no que se refere ao Componente de Conhecimento Específico – possuía boa capacidade de discriminar entre aqueles que dominam ou não o conteúdo.
O Índice de Facilidade variou de 0,14 a 0,87 e o de Discriminação, de 0,22 a 0,53.
Na Tabela 6.12, é apresentada a distribuição das questões levando -se em conta, simultaneamente, a classificação dos dois índices.
Dentre as questões que alcançaram os maiores Índices de Discriminação, uma delas foi classificada com Índice de Discriminação Muito bom, a de número 25.
Seu índice foi 0,53, e, quanto ao Índice de Facilidade, foi classificada como Médio, com 51 \% dos estudantes marcando a opção correta.
A Questão de número 32 foi classificada com Índice de Discriminação Muito bom ( 0,40 ), e, quanto ao Índice de Facilidade, foi classificada como Fácil, com 62 \% dos estudantes marcando a opção correta.
Já a Questão de número 21 foi a mais difícil entre as 27 questões específicas válidas, com baixo Índice de Facilidade, apenas 14 \% de acertos ( Muito Difícil ).
Essa questão apresentou poder discriminatório Médio, 0,22.
A título de exemplo das análises do comportamento das questões objetivas, no Gráfico 6.5, analisa -se a Questão 25 do Componente de Conhecimento Específico.
Essa questão foi considerada, pelas respostas dos estudantes avaliados na prova, como uma questão com nível de Facilidade Médio, 0,51, ou seja, 51 \% dos estudantes assinalaram, acertadamente, a opção C correspondente ao gabarito.
Como já comentado, seu Índice de Discriminação foi igual a 0,53, classificado como Muito bom, o maior valor de Discriminação.
Neste gráfico, \textbf{cada} uma das cinco curvas representa o percentual de respostas em determinada alternativa da Questão 25, em função do número de acertos dos estudantes nessa parte da prova, antes de possíveis eliminações de questões pelo critério Ponto-Bisserial.
A alternativa correta C, representada no gráfico pela curva em preto, foi escolhida em maiores proporções pelos estudantes com desempenho melhor nessa parte da prova.
Já as alternativas incorretas, também denominadas distratores, foram selecionadas, principalmente, por aqueles com notas mais baixas.
No caso de Música ( Licenciatura ), como se observa no eixo horizontal do Gráfico 6.5, nenhum estudante acertou mais do que 24 questões dentre as 27 questões objetivas válidas, considerando -se o gabarito original.
Observa -se que a soma não é 100 \%, por causa das questões não respondidas ou com mais de uma opção marcada.
Cerca de 50 \% daqueles com um acerto marcaram a opção A ( linha azul ), cerca de 25 \%, a opção B ( linha verde ) e cerca de 25 \%, a opção E ( linha vermelha ).
A proporção de estudantes que selecionaram a resposta correta C ( linha preta ) aumenta gradativamente, chegando a atingir 100 \% para 22, 23 e 24 acertos, enquanto a proporção dos que escolheram alternativas incorretas decai, a partir de nove acertos, em função do número de acertos nessa parte da prova.
Os gráficos relativos às demais Questões do Conhecimento Específico constam do Anexo I. 6.3 ANÁLISE DAS QUESTÕES DISCURSIVAS Esta seção apresenta estatísticas selecionadas e histogramas das Questões Discursivas de Formação Geral ( 6.3.1 ) e Conhecimento Específico ( 6.3.2 ).
São, também, apresentadas e comparadas as médias de subpopulações, caracterizadas por Grande Região.
Cumpre notar que nem todas as questões passam pelo crivo dos corretores.
As respostas depois de digitalizadas passam por um pré-processamento de identificação automático.
Algumas destas questões são definidas como um “ branco automático ” ( que não são enviadas para os professores para correção, a não ser que na segunda fase, a de inspeção visual, seja constatada uma resposta ).
O critério usado neste pré-processamento, é verificar para \textbf{cada} questão, a partir da quantidade de PIXELS, o que provavelmente, é uma resposta em branco.
A partir da separação destes documentos que foram considerados em branco, eles são enviados para uma verificação visual, por uma equipe de colaboradores, devidamente treinados.
Deste modo, quaisquer erros, deste pré-processamento, são encontrados, e as questões preenchidas, são enviadas para a correção dos professores.
Eventualmente, algumas respostas em branco, não são detectadas neste pré-processamento automático.
A presença de pequenos riscos, ou manchas, podem fazer o programa não identificar uma questão como um possível branco.
Estas são enviadas para a correção, mas são classificadas como “ branco ”, pelos professores quando da correção.
No caso da Área de Música ( Licenciatura ), foram identificados como brancos “ automáticos ”, 1.579 respostas de Questões Discursivas de Conhecimento Específico e 902 de Formação Geral.
Pelos professores corretores foram identificadas mais seis e sete questões em branco, respectivamente, para Conhecimento Específico e para Formação Geral.
Uma outra situação a ser mencionada é a ocorrências de atendimento a portadores de necessidades especiais.
O Anexo XI apresenta, para o exame como um todo e para a Área de Música ( Licenciatura ), por tipo de deficiência, o protocolo usado para permitir a correção das respostas dos alunos em \textbf{cada} situação. 6.3.1 COMPONENTE DE FORMAÇÃO GERAL As análises dos resultados de desempenho dos estudantes de Música ( Licenciatura ), nas duas Questões Discursivas relativas a Formação Geral, encontram -se na Tabela 6.13 e no Gráfico 6.6.
Aqui analisa -se a média das notas de \textbf{cada} Questão Discursiva de Formação Geral que leva em conta as correções realizadas por duas bancas, uma que analisa e pontua aspectos relativos ao conteúdo das respostas, seguindo o padrão de respostas divulgado pelo Inep ( ver Anexo IX ), outra que avalia o desempenho linguístico expresso pela redação das respostas, seguindo um padrão de respostas específico da Área de Língua Portuguesa ( ver Anexo IX ).
A nota de \textbf{cada} questão é composta de 80 \% da nota da correção de conteúdo e 20 \% da nota de desempenho linguístico.
Na Tabela 6.13, observa -se que a nota Média nesse conjunto de questões foi inferior à obtida nas objetivas.
Os estudantes de todo o Brasil obtiveram, em Formação Geral, Média 41,1 nas questões objetivas e 27,6 nas questões discursivas.
Pode -se notar, também, que o Desvio padrão nesse conjunto de questões foi menor do que o obtido nas objetivas : 19,9 nas questões objetivas e 17,4 nas questões discursivas.
A maior Média foi obtida na região Sul ( 31,1 ) e a menor, na região Nordeste ( 23,6 ).
A Mediana de todo o Brasil, neste componente, foi 27,5.
Nas regiões Norte, Nordeste e Centro-Oeste, a Mediana foi menor ( 24,0, 23,5 e 26,5, respectivamente ), e, nas demais regiões, foi maior que o valor para o Brasil : Sudeste ( 29,0 ) e Sul ( 32,3 ).
A nota Máxima ( 84,5 ) foi obtida na região Sudeste, sendo 64,5, a Máxima na região Norte ; 79,0, na Nordeste ; 74,5, na Sul ; e 70,5, na região Centro-Oeste. na Sul ; e 70,5, na região Centro-Oeste.
A nota Mínima ( 0,0 ) foi a mesma em todas as regiões do Brasil.
No Gráfico 6.6, \textbf{está} representada a distribuição das notas nas Questões Discursivas do Componente de Formação Geral.
A moda dessa distribuição ocorre no intervalo [ 20 ; 30 ], com frequência de 22,1 \% dos participantes, sendo que, no intervalo [ 0 ; 10 ], inclui -se, além da nota zero, a frequência de estudantes que deixaram esse tipo de questão em branco.
A distribuição é sempre decrescente a partir da moda.
Na sequência, os resultados verificados para \textbf{cada} uma das Questões Discursivas de Formação Geral serão apresentados, estabelecendo -se relações com os conteúdos abordados em \textbf{cada} uma delas.
Os comentários das Bancas de docentes corretores, a respeito do observado na correção das respostas dos estudantes, suas impressões e conclusões serão apresentados junto à análise de \textbf{cada} questão.
Cumpre esclarecer que, tendo em vista que as Questões Discursivas de Formação Geral são padronizadas, ou seja, constam de todas as provas, os comentários da Banca são os mesmos para todas as carreiras acadêmicas, sendo direcionados a todos os estudantes que participaram do Enade/2021.
A seguir, serão analisados o desempenho linguístico dos estudantes da Área de Música ( Licenciatura ) nas duas Questões Discursivas de Formação Geral do Enade/2021, apresentando -se os resultados obtidos e os comentários da banca de correção para \textbf{cada} questão. 6.3.1.1 ANÁLISE DE CONTEÚDO DA QUESTÃO DISCURSIVA 1 DO COMPONENTE DE FORMAÇÃO GERAL Os dados de Música ( Licenciatura ), obtidos a partir das respostas à questão 1, encontram -se na Tabela 6.14 e no Gráfico 6.7.
Nessa questão – de desempenho inferior ao da outra questão de Formação Geral –, os estudantes de todo o Brasil tiveram Média 19,8.
A maior Média para a Questão 1 foi obtida na região Sudeste ( 21,5 ), e a menor, na região Nordeste ( 16,4 ).
Quanto à variabilidade das notas, o Desvio padrão de todo o Brasil foi 19,9.
O menor Desvio padrão foi obtido na região Centro-Oeste ( 18,1 ), e o maior Desvio padrão foi obtido pela região Sudeste ( 20,8 ).
A Mediana do Brasil, como um todo, foi 15,0, sendo igual em todas as regiões, significando que pelo menos 50 \% dos participantes obtiveram essa nota ou nota menor.
A nota Máxima da Questão Discursiva 1 no Brasil foi 95,0, obtida por pelo menos um estudante da região Sudeste.
As demais notas Máxima foram : 75,0 ( Norte ), 85,0 ( Nordeste ), 80,0 ( Sul ) e 65,0 ( Centro-Oeste ).
A nota Mínima ( 0,0 ) foi a mesma em todas as regiões do Brasil.
No Gráfico 6.7, mostra -se a distribuição das notas na Questão Discursiva 1 do Componente de Formação Geral.
Observa -se que 13,0 \% dos participantes deixaram em branco a resposta a essa questão.
Somados à frequência dos que, tendo respondido, receberam nota zero, chega a 35,3 \% o contingente daqueles que zeraram questão.
Desconsiderando a proporção dos que zeraram a questão a moda seria o intervalo ( 10 ; 20 ], com 22,0 \% do total de participantes. 6.3.1.2 COMENTÁRIOS SOBRE A CORREÇÃO DE CONTEÚDO DAS RESPOSTAS À QUESTÃO DISCURSIVA 1 O enunciado da Questão 1 continha dois textos e um comando.
O primeiro texto trazia uma abordagem sobre as tensões entre cultura, arte e censura sob o recorte do inciso IX do Art. 5º da Constituição Federal de 1988, que diz ser livre a expressão da atividade intelectual, artística, científica e de comunicação, independentemente de censura ou licença.
O segundo texto apresentava o inciso IX do parágrafo 5º da Constituição Federal, em que é garantida a livre expressão da atividade artística, científica e de comunicação, independentemente de censura ou licença.
O comando solicitava que o respondente discorresse a respeito de dois aspectos.
Primeiramente, o estudante deveria ser capaz de dissertar sobre as relações entre os conceitos de arte, cultura e censura, à luz da ideia de liberdade artística garantida pela Constituição Federal de 1988.
O estudante também deveria ser capaz de apresentar duas ações educativas que pudessem contribuir “ para minimizar essas tensões ” ( tensões entre arte, cultura e censura ) e garantir a liberdade artística prevista pela lei.
A questão, portanto, \textbf{estava} de acordo com os temas previstos nas Diretrizes Curriculares para a Formação Geral : ética, democracia e cidadania ; cultura, arte e comunicação ; educação e desenvolvimento humano e social.
Além disso, o comando permitia avaliar o desenvolvimento das seguintes competências previstas nas Diretrizes Curriculares para a Formação Geral : planejar e elaborar projetos de ação e intervenção a partir da análise de necessidades em contextos diversos ; ler, produzir e interpretar textos com clareza e coerência ; identificar, compreender e analisar situações-problema utilizando pensamento lógico e sistêmico, estabelecendo relações de causalidade.
O enunciado da questão era claro, empregava linguagem adequada e continha informações suficientes para que os egressos construíssem seus argumentos.
Entretanto, apesar dessas características favoráveis, a temática era sofisticada, exigia reflexão e estabelecimento de relações a partir de textos concisos e densos.
Além disso, os conceitos de arte e cultura expressos pelo Texto I excluíam inúmeras outras definições reconhecidas em diferentes áreas das ciências humanas, em alguns casos divergentes daquelas apresentadas.
Os conhecimentos e o nível de profundidade exigidos pelo padrão de resposta eram adequados para egressos de cursos superiores.
Deve -se destacar que o padrão admitia um espectro amplo para a construção da argumentação das respostas, viabilizando a aceitação de outras interpretações dos conceitos de arte e cultura, incluindo possíveis contestações em relação às ideias apresentados no Texto I.
A questão pode ser considerada difícil, já que a primeira parte exigia uma resposta que articulasse quatro conceitos complexos – arte, cultura, censura e liberdade artística –, e a segunda parte exigia o conhecimento do que seja uma ação educativa, além da proposição de ações que objetivamente contribuíssem para minimizar as tensões entre arte, cultura e censura, além de garantir a liberdade artística.
Essas observações, aliadas ao fato do espaço e do tempo limitados para elaboração de uma resposta tão extensa, elevaram consideravelmente o nível de dificuldade da questão.
A correção confirmou a dificuldade dos estudantes para abordar a temática, evidenciou grande diversidade na compreensão dos conceitos envolvidos e falta de clareza do que seja uma ação educativa.
Nas respostas à primeira parte do comando, os estudantes raramente abordaram os quatro conceitos que deveriam \textbf{estar} presentes na resposta.
A maioria se limitou a relacionar arte e cultura com comentários que repetiam trechos dos enunciados, como arte e cultura"andam de mãos dadas".
Assim, apesar de o Texto I sugerir o entendimento de que a cultura é uma esfera reguladora do comportamento e de que a arte é o elemento de desequilíbrio, de contestação, a quase totalidade das respostas não fez uso dessa interpretação, tratando arte e cultura quase como sinônimos.
Apesar de raras, houve respostas nas quais os estudantes defenderam a existência de múltiplas culturas, muitas delas não hegemônicas, tão suscetíveis aos dispositivos de controle quanto a arte.
Houve também respostas que colocavam em questão os critérios, muitas vezes elitistas e eurocêntricos, utilizados para produzir hierarquização entre as artes, chegando a inviabilizar o reconhecimento de outras tantas manifestações artísticas.
Foi pequeno o grupo que, além de relacionar arte e cultura, foi capaz de articulá -las com o conceito de censura trazendo -a, quase sempre, como herança do período militar.
Entretanto, nesse grupo, outras abordagens de censura foram mapeadas no decorrer das correções : cancelamento em redes sociais ; contrainformação ou manipulação dos meios de comunicação social ; além de outros recursos censórios mais indiretos, como os processos judiciais, a suspensão de financiamentos públicos a projetos culturais e a autocensura.
Em relação à segunda parte do comando, era esperado que as/os estudantes egressas/os citassem duas ações educativas sob o recorte de que minimizassem tensões observadas entre arte, cultura e censura e garantissem a liberdade artística prevista pela lei.
Na maioria das respostas, não se verificou essa relação entre a proposição e seu objetivo/resultado ; muitas ações foram apenas citadas sem explicações.
Esse fato reduziu bastante a quantidade de respostas com pontuação integral no referido item.
Importante ressaltar que o padrão de resposta indicou exemplos de ações educativas no campo da educação não formal.
Entretanto, por não ser algo explícito no comando, foram também aceitas ações referentes ao campo da educação formal.
O mais recorrente nas respostas foi a indicação de ações educativas em sala de aula e visitas a museus e a exposições, mediadas pelo próprio professor da turma.
Entretanto, houve respostas em que os estudantes ampliaram tanto o espaço quanto o grupo considerado nas ações educacionais, propondo campanhas para disseminar conhecimento sobre arte e/ou cultura para plateia maior do que a de estudantes.
De modo geral, a fluidez do texto das respostas e a consistência das ideias articuladas ficou aquém do padrão desejado para estudantes egressos do ensino superior.
O maior volume de respostas mostrou que os estudantes valorizaram mais a apresentação de duas ações educativas, muitas vezes de forma incompleta, apenas citando uma ação proposta e deixando de explicar o desdobramento solicitado no comando.
Além disso, houve respostas em que sugeriram, de forma isolada, ações do tipo apresentações, shows, exposições, sem qualquer \textbf{detalhamento}.
Entretanto, uma ação educativa com as finalidades solicitadas precisaria ser minimamente descrita para que fosse possível identificar que a ação pressupunha a presença de um mediador ( professor / arte educador / museólogo ) e de receptores participativos.
A proposição de uma ação educativa também foi confundida pelos estudantes com a proposta de implantação de equipamentos de infraestruturas ou projeto governamental, tais como : escolas, centros educacionais, incentivo ao ensino das artes em áreas mais necessitadas e introdução de programas de ensino de artes nos currículos escolares.
O contexto da arte e/ou da cultura e seus confrontos com a censura foram considerados, muitas vezes, sem o estabelecimento de relações, mas restringindo -se a aspectos descritivos de como o respondente compreendia \textbf{cada} um desses elementos.
Há de se considerar, adicionalmente, que em uma parte efetiva das respostas os estudantes defendiam, indevidamente, a censura como um mecanismo legitimamente regulador dos chamados"excessos da arte", revelando confusão entre os limites legais estabelecidos em lei ( que tipificam discursos de ódio, pedofilia, apologia às drogas, etc. ) e os julgamentos pessoais, coletivos ou institucionais que desencadeiam atitudes censoras.
A abrangência proporcionada pelo tema e pelo comando refletiu -se nas respostas, principalmente nos exemplos das ações educativas apresentados pelos estudantes.
Já na reflexão sobre arte, cultura e censura não houve tanta diversidade – a discussão quase sempre permanecia no senso comum, com menção à influência da época militar ou associada 6.3.1.3 ANÁLISE DE CONTEÚDO DA QUESTÃO DISCURSIVA 2 DO COMPONENTE DE FORMAÇÃO GERAL A Tabela 6.15 mostra que o desempenho médio dos estudantes, na Questão Discursiva 2 ( Média 24,8 ), foi maior que o obtido na Questão Discursiva 1 ( Média 19,8 ).
A região Nordeste foi aquela cuja Média foi menor ( 20,3 ), e a maior Média ocorreu na região Sul ( 30,1 ).
Quanto à variabilidade das notas, o Desvio padrão de todo o Brasil foi 22,8, um pouco maior ao obtido na Questão Discursiva 1 ( 19,9 ).
O maior desvio nessa questão foi obtido na região Sul ( 23,9 ), enquanto o menor foi obtido na região Norte ( 21,2 ).
A Mediana de todo o Brasil foi 25,0, a mesma de das regiões Sudeste e Sul Nas regiões Norte e Nordeste a Mediana foi 15,0 e na Centro-Oeste foi 10,0.
A nota Máxima ( 100,0 ) foi obtida nas regiões Sudeste e Sul, sendo 95,0, a nota Máxima nas regiões Norte e Nordeste, e 80,0, na região Centro-Oeste.
A nota Mínima ( 0,0 ) foi a mesma em todas as regiões do Brasil. à discussão atual sobre limitações à liberdade de expressão.
Houve grande incidência de respostas em que o elemento censura foi confundido com a noção de classificação indicativa.
O desempenho ficou aquém do esperado.
Na maioria das respostas, os alunos ficaram com notas baixas, e foi grande a incidência de zeros.
Cabe ressaltar também a grande quantidade de respostas em branco e as que receberam nota zero devido a anulações por defesa de atos de censura, na contramão do que estabelecem a Constituição Federal e os Direitos Humanos.
A fuga ao que foi solicitado pelo comando da questão foi marcante, talvez evidenciando dificuldade de interpretação ou distanciamento em relação ao tema.
Outra lacuna muito observada foi a dificuldade de se expressar em textos fluentes e articulados.
A fragmentação e a falta de profundidade da argumentação foram uma constante.
Além disso, chamou a atenção a predominância do lugar comum na construção das respostas.
Tratando -se de um universo de formandos de nível superior, a superficialidade da reflexão em \textbf{torno} da arte e da censura e a dificuldade para incluir a Constituição Federal na discussão foram surpreendentes. 6.3.1.3 ANÁLISE DE CONTEÚDO DA QUESTÃO DISCURSIVA 2 DO COMPONENTE DE FORMAÇÃO GERAL A Tabela 6.15 mostra que o desempenho médio dos estudantes, na Questão Discursiva 2 ( Média 24,8 ), foi maior que o obtido na Questão Discursiva 1 ( Média 19,8 ).
A região Nordeste foi aquela cuja Média foi menor ( 20,3 ), e a maior Média ocorreu na região Sul ( 30,1 ).
Quanto à variabilidade das notas, o Desvio padrão de todo o Brasil foi 22,8, um pouco maior ao obtido na Questão Discursiva 1 ( 19,9 ).
O maior desvio nessa questão foi obtido na região Sul ( 23,9 ), enquanto o menor foi obtido na região Norte ( 21,2 ).
A Mediana de todo o Brasil foi 25,0, a mesma de das regiões Sudeste e Sul Nas regiões Norte e Nordeste a Mediana foi 15,0 e na Centro-Oeste foi 10,0.
A nota Máxima ( 100,0 ) foi obtida nas regiões Sudeste e Sul, sendo 95,0, a nota Máxima nas regiões Norte e Nordeste, e 80,0, na região Centro-Oeste.
A nota Mínima ( 0,0 ) foi a mesma em todas as regiões do Brasil.
No Gráfico 6.8, mostra -se a distribuição das notas na Questão Discursiva 2 do Componente de Formação Geral.
Observa -se que 16,3 \% dos estudantes deixaram a questão em branco.
Chega a 26,3 \% \% o percentual de notas zero agregando aos brancos a frequência daqueles que, tendo respondido, não pontuaram.
A moda é o intervalo ( 20 ; 30 ], com frequência de 24,4 \%.
Nota -se, ainda, que as notas ficaram mais dispersas em comparação às da Questão Discursiva de número 1, o que pode ser constatado, também, pela comparação do Desvio padrão das notas da Questão Discursiva 2 ( 22,8 ) e o das notas da Questão Discursiva 1 ( 19,9 ). 6.3.1.4 COMENTÁRIOS SOBRE A CORREÇÃO DE CONTEÚDO DAS RESPOSTAS À QUESTÃO DISCURSIVA 2 O enunciado da segunda Questão Discursiva de Formação Geral era composto por dois textos motivadores e um comando com dois itens.
O primeiro texto apresentava uma definição de cidade inteligente, pelo uso da tecnologia nos diversos setores da administração e seu compromisso com o meio ambiente.
O segundo texto ampliava a definição anterior explorando aspectos da gestão pública.
O enunciado era claro, empregava linguagem adequada e continha informações suficientes para que os egressos construíssem seus argumentos.
O comando solicitava, primeiramente, que o estudante explicasse de que modo as cidades inteligentes podem contribuir para a melhoria das questões relacionadas ao desenvolvimento sustentável.
Em atendimento ao segundo item do comando, o estudante deveria apresentar uma proposta de intervenção urbana que gerasse impacto social e contribuísse para a melhoria da vida em comunidade.
Os dois itens do comando exigiam articulação de ideias e exposições com desdobramentos do porquê da resposta, o que tornou a questão complexa.
A questão abordava, portanto, temas previstos nas Diretrizes Curriculares para a Formação Geral : ciência, tecnologia e inovação ; biodiversidade, sustentabilidade e intervenção humana e cidades, habitação e qualidade de vida.
Além disso, avaliava o desenvolvimento de competências também previstas nas Diretrizes Curriculares para a Formação Geral : planejar e elaborar projetos de ação e intervenção a partir da análise de necessidades em contextos diversos ; ler, produzir e interpretar textos com clareza e coerência ; identificar, compreender e analisar situações-problema utilizando pensamento lógico e sistêmico, estabelecendo relações de causalidade.
Os conhecimentos e o nível de profundidade exigidos foram considerados adequados para egressos de cursos superiores pelos componentes da banca de correção.
Os textos tratavam de temas atuais e de amplo debate na mídia, tais como : cidade inteligente, meio ambiente, desenvolvimento sustentável, intervenção urbana e impacto social.
No entanto, o nível de dificuldade foi considerado elevado, já que tais conceitos admitem interpretações diversas nesse universo particular de egressos, o que criou, em consequência, um espectro diversificado de respostas e tipos de argumentação.
Além disso, era necessário conhecer o conceito de intervenção urbana para propor uma ação deste tipo que gerasse impacto social.
O padrão de resposta \textbf{estava} redigido de forma clara e, para o item ‘ a ’ do comando, apresentava um exemplo simples e informativo de como a tecnologia pode ser usada para a redução do consumo de energia e emissão de CO2.
Para o item ‘ b ’, o padrão apresentava exemplos de propostas de intervenção de forma mais diversificada do que a resposta padrão do item ‘ a ’, levantando várias opções corretas.
A abordagem adotada pelos estudantes foi quase sempre a da exemplificação, o mesmo recurso adotado no padrão de resposta.
Assim, pode -se dizer que as respostas dos estudantes se mantiveram, quase sempre, em conformidade com o contexto do padrão, mas com uma variedade muito grande de temas e ambientes relacionados ao desenvolvimento sustentável.
Dentre as ações vinculadas às cidades inteligentes, as ocorrências mais significativas foram as relacionadas à poluição, ao desmatamento, aos transportes coletivos, à segurança e ao desenvolvimento de tecnologia de modo a contribuir para a melhoria das dificuldades enfrentadas atualmente nas cidades.
Em muitos casos, o estudante articulou a sua resposta com a realidade vivenciada em determinadas cidades, comparando, inclusive, cidades que são consideradas inteligentes com realidades locais.
Também foram encontradas respostas em que os estudantes apresentavam propostas e soluções para realidades específicas das cidades onde o respondente reside, contextualizando o uso da tecnologia com ações fundamentais que faltam em algumas cidades brasileiras, como saneamento básico, água não poluída para uso doméstico, destinação do lixo.
As respostas incompletas ocorreram, muitas vezes, em função de o aluno apenas apresentar ações ou características das cidades inteligentes sem articular com o desenvolvimento sustentável ou com a intervenção urbana.
Em muitas respostas, os alunos apenas repetiram ideias do enunciado ; em outras, demonstraram a não compreensão sobre o que \textbf{estava} sendo solicitado pelo comando, deixando de argumentar sobre uma proposta ou ação mencionada.
Nas respostas, os estudantes se concentraram no uso da tecnologia para monitorar o desmatamento e a violência ; em hortas comunitárias para melhorar a qualidade da alimentação das pessoas ; em investimento em locais públicos para a realização de exercícios ; em maior controle do uso de combustíveis não fósseis para evitar a poluição.
As propostas de intervenção fixaram -se bastante em aplicativos, solução sugerida pelo ambiente atual de desenvolvimento e das facilidades proporcionadas pela Internet.
Marcante foi a diversidade de funções encontradas nas propostas de sensores e aplicativos, desde marcador de horário de passagem de ônibus até controlador de desperdício de algum bem, como água, por exemplo.
O resultado ficou abaixo do esperado.
A quantidade expressiva de zeros talvez tenha sido consequência de desinteresse mais do que por dificuldades para interpretar e responder ao comando.
Essas ficaram mais evidentes nas propostas de intervenção urbana, muitas vezes incompletas por sequer explicitarem o benefício delas advindo.
As maiores fragilidades reveladas pelas respostas dos estudantes, apontadas pelos corretores, foram a interpretação do contexto e do comando da questão, além de uma grande dificuldade para se expressarem de forma fluente e concisa.
As respostas com geração de ideias próprias ocorreram em escala bem menor do que aquelas em que os estudantes usaram do recurso da paráfrase pobre sobre o texto do enunciado, revelando dificuldade para criar e articular ideias a respeito dos temas em foco.
A linguagem utilizada nas respostas ficou, muitas vezes, aquém do padrão desejado para estudantes egressos do ensino superior.
A clareza das ideias apresentadas e a articulação do conteúdo com o comando da questão raramente foram encontradas nas respostas no nível esperado.
A dificuldade mais evidente dos estudantes foi a de estabelecer relações e encadear ideias para completar um raciocínio iniciado ; foram muitas as respostas fragmentadas. 6.3.1.5 CONSIDERAÇÕES FINAIS SOBRE A CORREÇÃO DE CONTEÚDO DAS QUESTÕES DISCURSIVAS DO COMPONENTE DE FORMAÇÃO GERAL As Questões Discursivas da Área de Formação Geral abordaram temas atuais e proporcionaram a avaliação de competências importantes dentre as previstas pelas suas Diretrizes Curriculares.
A análise das respostas revelou, principalmente, dificuldade dos estudantes na interpretação dos enunciados – talvez pela linguagem dos textos motivadores ser mais complexa do que a do cotidiano – e na compreensão dos comandos das questões – que pode ser fruto de desatenção ou de desconhecimento mais profundo do contexto.
Esses fatos parecem ter sido os principais responsáveis pelo resultado aquém do esperado em termos de avaliação quantitativa.
Mesmo os estudantes que responderam em conformidade com o padrão demonstraram um universo de argumentação muito limitado, pobre em vocabulário, com o recurso frequente de procurar nos textos introdutórios os elementos para a construção de suas frases.
Isso criou uma peculiaridade interessante : grupos de respostas muito parecidas.
Encadear e relacionar argumentos também foi uma deficiência percebida.
As razões eram apresentadas de forma confusa, sem fluência e sem constituir unidade.
O conteúdo das respostas aponta uma preocupação com o viés social dos problemas abordados ; são bons exemplos as propostas bastante diversificadas de melhoria de vida para a parcela da população menos favorecida.
Por outro lado, traz indícios de deficiências consideráveis na separação conceitual de cultura e arte, na compreensão mais ampla do conceito de censura e na capacidade de separar ações educativas de culturais ; em algumas respostas, também se evidenciaram conceitos equivocados sobre tecnologia, com exemplos de concretização muito difícil física ou financeiramente.
Houve alguns casos em que os conceitos de sustentabilidade e de meio ambiente se confundiam em lugar de se complementarem.
O recorte do grupo de formandos avaliado, em resumo, é de dificuldade de interpretação e produção de textos, vocabulário restrito ao praticado cotidianamente, assunção de fraseologia de senso comum em lugar do conteúdo apreendido formalmente e falta de profundidade em matérias importantes para a formação do cidadão como cultura, arte, educação, liberdade de expressão, tecnologia, cidade inteligente e desenvolvimento sustentável. 6.3.1.6 ANÁLISE DE LÍNGUA PORTUGUESA DAS QUESTÕES DISCURSIVAS DO COMPONENTE DE FORMAÇÃO GERAL Os dados de Música ( Licenciatura ), obtidos a partir das respostas às Questões Discursivas do Componente de Formação Geral, no que tange à Língua Portuguesa, encontram -se na Tabela 6.16 e no Gráfico 6.9.
Nesse aspecto, os estudantes, de todo o Brasil, obtiveram Média 48,6.
A maior Média em relação à Língua Portuguesa foi obtida na região Sul ( 53,2 ), e a menor, na região Norte ( 43,0 ).
Quanto à variabilidade das notas, o Desvio padrão de todo o Brasil foi 23,4.
O menor Desvio padrão foi obtido na região Sul ( 20,0 ), e o maior Desvio padrão foi obtido na região Centro-Oeste ( 26,2 ).
A Mediana das notas de Língua Portuguesa foi 55,0 para o Brasil, como um todo.
Nas regiões Sudeste e Sul, a Mediana foi 57,5, nas regiões Nordeste e Centro-Oeste, 52,5 ; e na região Norte, a Mediana foi 47,5.
A nota Máxima para todo o Brasil foi 90,0, com, pelo menos, um estudante tirando essa nota nas regiões Sudeste e Sul.
Nas demais regiões a nota Máxima foi 85,0 ( Norte e Nordeste ) e 82,5 ( Centro-Oeste ).
A nota Mínima foi zero em todas as regiões do país.
No Gráfico 6.9, mostra -se a distribuição das notas de Língua Portuguesa do Componente de Formação Geral.
Observa -se que a maior frequência ( 24,4 \% ) corresponde à dos estudantes que obtiveram nota no intervalo ( 50 ; 60 ].
Destacam -se, também, os estudantes que deixaram as duas questões discursivas de Formação Geral em branco, 10,6 \% do total. 6.3.1.7 COMENTÁRIOS SOBRE A CORREÇÃO DE LÍNGUA PORTUGUESA DAS RESPOSTAS ÀS QUESTÕES DISCURSIVAS A avaliação de desempenho linguístico De acordo com a portaria nº 386 do MEC/INEP, de 23 de agosto de 2021, o Enade deve avaliar se o estudante desenvolveu, no processo de formação, competências para : compreender as linguagens e respectivas variações ( item V ) ; ler, produzir e interpretar textos com clareza e coerência ( item VI ) ; analisar e interpretar representações verbais, não verbais de fenômenos diversos ( item VII ) ; identificar diferentes representações de um mesmo significado ( item VIII ) ; formular e articular argumentos e contra-argumentos consistentes em situações sociocomunicativas ( item IX ).
Essas competências são essenciais à formação de nível superior, porque dizem respeito ao domínio da linguagem, tanto no processo de leitura quanto na articulação e na organização de ideias para a produção textual.
Dessa forma, as duas questões discursivas de formação geral oferecem oportunidade para avaliar se o participante consegue desenvolver um tema, recorrendo a argumentos que explicitem sua opinião e elaborando propostas de intervenção social adequadas a um exercício consciente da cidadania, conforme determina a referida portaria : “ As questões discursivas do componente de Formação Geral avaliam aspectos como clareza, coerência, coesão, estratégias argumentativas, vocabulário e gramática adequados à norma-padrão da Língua Portuguesa. ” ( Art.2 §2 da referida portaria nº 386 ) Para avaliar esses aspectos, portanto, é necessário identificar, nas respostas discursivas de Formação Geral, as diferentes habilidades necessárias ao bom desempenho linguístico na situação de formalidade exigida : a organização das ideias, tendo em vista a construção de um texto coerente e estruturado ; o domínio das convenções ortográficas e das convenções morfossintáticas que constituem o arcabouço da gramática da norma-padrão da Língua Portuguesa ; e o domínio de um vocabulário amplo e diversificado, que permita uma comunicação clara e eficiente.
A Questão Discursiva 1 desenvolveu o tema da relação entre arte, cultura e censura, tomando como base dois fragmentos de textos.
O enunciado \textbf{estava} claro ao encaminhar a solicitação de resposta, evidenciando a exigência de elaboração de um texto dissertativo, por meio do verbo"discorrer", e de apresentação de duas ações educativas, articuladas com o ponto de vista desenvolvido na parte argumentativa.
A contraposição dos dois textos favoreceu a formulação e a articulação de argumentos e contra-argumentos consistentes, em defesa do ponto de vista a ser defendido pelo participante.
O enunciado permitia, inclusive, a ampliação da abordagem com base na análise de situações concretas que evidenciassem essa articulação.
A Questão Discursiva 2 desenvolveu o conceito de"cidade inteligente", tomando como base dois fragmentos de estudos que descreviam as suas características.
Com base nesses dois fragmentos, o comando solicitou do participante :"explique a contribuição das cidades inteligentes para o desenvolvimento sustentável"e apresente uma proposta de intervenção urbana que contribua para a melhoria da vida em comunidade.
O comando “ explicar ” demandava a elaboração de um texto de base expositiva, sem desenvolvimento de argumentação favorável ou contrária à configuração de uma cidade inteligente e de seus benefícios para a sociedade.
Em relação à segunda parte, o comando exigia a elaboração de uma proposta de intervenção urbana, que deveria conter informações sobre os agentes da intervenção, os recursos aplicados e o \textbf{detalhamento} dessa ação.
A segunda questão, portanto, se diferenciava da primeira, por solicitarem estruturas textuais distintas : a Questão 1 se consolidava com base em um modo de organização argumentativo e a Questão 2 se consolidava com base em um modo de organização expositivo.
Esses dois estilos textuais ofereceram oportunidade de avaliar se os participantes \textbf{estavam} habilitados a desenvolverem textos com esses dois distintos objetivos comunicativos, o que seria esperado ao final do curso de graduação e necessário à sua vida profissional.
A análise do desempenho linguístico esperado de um estudante que \textbf{está} concluindo seu curso superior apoiou -se no princípio de que ele deveria utilizar conhecimentos sobre os temas e estruturar seus textos de acordo com as características do registro formal adequado à situação comunicativa – avaliação de conhecimentos ao final de um curso superior.
Essa configuração determina exigências quanto à adequação da seleção vocabular, ao desenvolvimento do conteúdo, à estruturação sintática dos períodos, à organização lógica das ideias, à utilização de procedimentos de encadeamento textual e de referenciação, à obediência às exigências morfossintáticas próprias da modalidade escrita da norma-padrão, ao respeito às regras ortográficas de acentuação gráfica e de grafia das palavras.
O padrão de resposta utilizado na avaliação A partir desses preceitos, foram considerados, no padrão de resposta, as seguintes competências : ( 1 ) estruturação textual ; ( 2 ) respeito às convenções ortográficas da norma-padrão da Língua Portuguesa ; ( 3 ) domínio dos diferentes aspectos morfossintáticos próprios da modalidade escrita formal da norma-padrão da Língua Portuguesa ; e ( 4 ) seleção vocabular adequada à modalidade escrita formal da Língua Portuguesa, exigida pela situação comunicativa.
A seguir, os aspectos avaliados por \textbf{cada} uma destas competências, são apresentados. ( 1 ) Estruturação textual condizente com o gênero textual solicitado e o modo de organização textual adequado ao gênero ( argumentativo ou expositivo ).
Esta competência envolve : - estruturação sintática condizente com o padrão da modalidade escrita formal da Língua Portuguesa de modo a garantir a clareza e a coerência necessárias ; - distribuição do conteúdo do texto em parágrafos, de modo a garantir a sua organização temática ; - distribuição das ideias em períodos que mantenham uma estruturação sintática condizente com a modalidade formal da língua ( evitando -se frases fragmentadas, sequência justaposta de ideias sem encaixamentos sintáticos e presença de oração subordinada sem oração principal que reflitam hábitos da oralidade e comprometam a estrutura lógico-gramatical do texto ) ; - utilização de operadores discursivos que contribuam para a progressão temática do texto, estabelecendo relações lógicas entre as ideias apresentadas, tanto do ponto de vista intrafrasal, como interfrasal : preposições, conjunções, pronomes relativos, advérbios e locuções adverbiais de tempo e lugar ; - utilização de procedimentos de referenciação lexical e pronominal que permitam a retomada de referentes textuais, como pronomes, advérbios e artigos ; além de elementos lexicais, como sinônimos, hiperônimos, hipônimos, expressões metafóricas e expressões resumitivas ) ; - emprego adequado dos pronomes relativos “ cujo(a) ” e “ onde ”, responsáveis por encaixamentos sintáticos no período composto ; - utilização adequada de sinais de pontuação que contribuam para a organização lógica da frase e do texto ; - inteligibilidade relacionada ao atendimento das exigências de estruturação textual.
A expectativa, portanto, era de que o participante recorresse a procedimentos linguístico-discursivos para organizar seu texto, permitindo o encadeamento lógico entre suas partes de forma a garantir a progressão e a coerência textuais, com estruturação que não reproduzisse hábitos da oralidade e sim atendesse à norma-padrão da Língua Portuguesa. ( 2 ) Respeito às convenções ortográficas da norma-padrão da Língua Portuguesa.
Esta competência envolve o respeito às convenções estabelecidas pela legislação em vigor e consubstanciadas no Vocabulário Ortográfico da Língua Portuguesa, editado pela Academia Brasileira de Letras e avaliam -se : - domínio das regras de acentuação gráfica ; - domínio da grafia padrão das palavras ( com ausência de abreviaturas próprias da linguagem da internet ).
Assim, esperava -se que o participante : grafasse corretamente as palavras ; respeitasse as regras de acentuação gráfica ; empregasse maiúsculas em início de frase, em nomes próprios de pessoas, lugares ou instituições ; grafasse as siglas com letras maiúsculas ou apenas com a primeira letra maiúscula quando formarem uma palavra ; evitasse abreviações como p/, vc, tb, pq, \textbf{tá}, né, usadas muitas vezes em escrita informal e na internet ; obedecesse às regras de separação de sílabas no final da linha ; atendesse aos padrões de segmentação vocabular da norma-padrão. ( 3 ) Domínio dos diferentes aspectos morfossintáticos próprios da modalidade escrita formal da norma-padrão da Língua Portuguesa.
Esta competência envolve o domínio das convenções próprias da norma padrão da Língua Portuguesa relativas aos seguintes aspectos : concordância nominal ; concordância verbal ; regência nominal ; regência verbal ; flexão nominal ; flexão verbal ; correlação entre os tempos verbais ; colocação pronominal.
Assim, esperava -se que o participante : - flexionasse o verbo para estabelecer concordância de número com o sujeito da frase ; - flexionasse o artigo, o adjetivo e o pronome para concordar em número e em gênero com o substantivo a que \textbf{cada} um se refere ; - observasse a regência nominal e a verbal, utilizando a preposição adequada depois de um substantivo, um verbo ou um adjetivo ; - empregasse adequadamente o acento grave indicador da crase que expresse a contração entre uma preposição e um artigo, de acordo com as exigências de regência nominal ou verbal ; - obedecesse às regras de colocação pronominal ( próclise e ênclise ), distintas dos hábitos da oralidade ou da escrita informal ; - flexionasse adequadamente verbos, substantivos, adjetivos e pronomes, no que diz respeito à expressão das categorias gramaticais ; - flexionasse os verbos para expressar a correlação de modo e de tempo nas estruturas subordinadas.
Com base nesses critérios, alguns desvios ortográficos foram considerados como desvios de caráter morfossintático porque interferem na classificação dos vocábulos de base morfológica e funcional : - eliminação da marca de infinitivo ( -r- ) e substituição por acento agudo ou ausência total de marca do infinitivo, causando ambiguidade entre distintas classes de palavras ; - confusão entre as terminações “ ão ” e “ am ”, utilizadas para marcar diferentes tempos verbais ; - emprego indevido da forma verbal de"haver"\&lt ; “ há”\&gt ; no lugar da preposição ou do artigo \&lt ; “ a”\&gt ;, e vice-versa ; - emprego de hífen para separar pronome átono – tanto uso indevido quanto omissão ( ex : “ esperasse ”, em lugar de “ espera -se ” ; “ falar -mos ”, no lugar de “ falarmos ” ) ; - emprego de “ esta ” no lugar de “ \textbf{está} ”, provocando ambiguidade devido à confluência gráfica entre o pronome demonstrativo e a forma do verbo"\textbf{estar}", pela ausência da acentuação gráfica ; - emprego da conjunção aditiva"e"no lugar da forma verbal “ é ”, provocando ambiguidade devido à confluência gráfica por ausência da acentuação gráfica ; - emprego do advérbio “ mais ” no lugar da conjunção adversativa “ mas ”, reproduzindo confluência fônica da oralidade ; - registro da terceira pessoa do plural dos verbos “ ter ” e “ vir ” sem o acento circunflexo, causando desvio de concordância verbal ; - emprego inadequado do relativo “ onde ” deslocado do seu sentido locativo, reproduzindo tendência da oralidade ; - emprego do pronome relativo “ que ” sem a necessária preposição ( desvio de regência ) ; - ausência do acento grave obrigatório para marcar a crase. ( 4 ) Seleção vocabular adequada à modalidade escrita formal da Língua Portuguesa, exigida pela situação comunicativa.
Esta competência envolve : - precisão na seleção/utilização do vocabulário relacionado à temática solicitada pela questão ; - diversificação e adequação vocabular ; - coerência na combinação vocabular ; - ausência de marcas da oralidade, como termos de sentido muito genérico ( “ coisa ”, “ negócio ”, “ você ” ) e termos de registros mais informais ( como gírias, jargões, frases feitas, ditados populares, termos regionais ).
Assim, esperava -se que o participante : respeitasse a adequação vocabular, não empregando gírias ou expressões coloquiais ; apresentasse repertório vocabular diversificado e preciso ; evitasse repetição desnecessária de palavras ; utilizasse um vocabulário mais formal, como solicitado por um texto dissertativo.
Essas quatro competências foram reunidas em três itens de avaliação, atendendo aos aspectos explicitados anteriormente : I - Aspectos Ortográficos ( 20 \% ) : domínio das convenções ortográficas : grafia de vogais e consoantes, uso de maiúsculas e minúsculas, emprego do hífen e da acentuação gráfica, segmentação de palavras ; II - Aspectos Textuais ( 40 \% ) : domínio dos procedimentos de estruturação textual : organização interna dos períodos, emprego de conectores para a articulação lógica entre os períodos e entre os parágrafos, emprego de marcas de referenciação lexical e pronominal ; utilização dos sinais de pontuação que contribuem para a organização lógica da frase ; III - Aspectos Morfossintáticos e Vocabulares ( 40 \% ) : domínio das regras de caráter morfossintático estabelecidas como modelares do ponto de vista da modalidade escrita formal da norma-padrão da Língua Portuguesa : concordância nominal e verbal, regência nominal e verbal, colocação pronominal, flexão nominal e verbal, correlação entre tempos e modos verbais, ausência de marcas de oralidade.
Outra competência combinada a esta na análise das respostas é a seleção vocabular adequada à modalidade escrita formal da Língua Portuguesa, tendo em vista a intersecção entre as duas do ponto de vista das exigências do registro formal da modalidade escrita da norma-padrão.
Análise das respostas dadas pelos estudantes A Questão 1 propiciou a oportunidade de maior desenvolvimento do tema solicitado, gerando consequências na estruturação textual das respostas, já que os textos foram mais opinativos, mais elaborados e fluentes.
Já a Questão 2 teve um comportamento distinto, devido ao seu caráter mais expositivo, os textos foram mais curtos e o vocabulário mais limitado.
Observaram -se, nas duas questões, alguns casos de comprometimento de autoria, com textos muito próximos a paráfrases dos textos motivadores e sem marcas de repertório sociocultural produtivo.
Apresentam -se, a seguir, algumas observações sobre o desempenho dos participantes : I - Aspectos Ortográficos : o desempenho dos participantes revelou uma diferença muito grande nos dois aspectos analisados, como já foi observado nos relatórios dos anos anteriores : baixo índice de desvios da grafia padrão e grande índice de desvios de acentuação.
Em vários casos, ocorre ausência quase completa de acentuação gráfica.
Vale ressaltar que muitos participantes apresentaram uma caligrafia quase ininteligível.
Ainda nesse aspecto, ocorreu a omissão inadequada de cedilha, de til e do traço do T em algumas respostas.
Houve, inclusive, respostas com ausência total de cedilha ou til.
Os casos mais sistemáticos de eliminação do acento indicador da sílaba tônica foram : palavras proparoxítonas ( “ artistica ”, “ ludico ”, “ publico ”,"musica") ; palavras paroxítonas terminadas em ditongo crescente ( “ necessario ”,"historia","propria", “ vivencias ” ) ; palavras oxítonas ( “ tambem ”,"populacao").
Por outro lado, destaca -se o uso indevido do acento gráfico em palavras como “ propíciar ” “ criânça ”, por exemplo.
Quanto aos desvios de grafia, observam -se os seguintes exemplos : emprego inadequado de maiúscula : “ a Liberdade artística ” ; ou de minúscula :"constituição"; confusão entre"ss","sc","c"e"ç": “ condenassão ”, “ nessessario", “ repreiensão ”, “ conciencia ” ; eliminação de grafemas : “ essescia ”,"questõe"; - desvios de segmentação : “ ser vidor", “ igual dade ”, “ possi bilidades"; troca de grafema : “ convença ” ( por “ conversa ” ), “ comun ”.
II - Aspectos Textuais : a estruturação textual foi a competência que apresentou o maior número de desvios, sendo responsável, em alguns textos, pela falta de clareza e pela desorganização.
No conjunto dos textos analisados, observaram -se diferentes níveis de domínio dos recursos de estruturação e organização textual.
No entanto, os participantes podem ser agrupados em três faixas descritas a seguir : - os que possuem grande domínio dos recursos de estruturação textual, com produção de textos que atendem às exigências de clareza e coerência, com progressão temática bem direcionada e utilização produtiva de operadores argumentativos ; - os que apresentam um domínio regular dos recursos de estruturação textual, com desvios na organização do texto e ausência de paragrafação, além de uso insuficiente de operadores argumentativos e processos de encaixamento sintático ; - os que apresentam baixo domínio dos recursos de estruturação textual, com muitos desvios que retomam hábitos da oralidade e que comprometem a coerência do texto : ocorrência de sequência justaposta de ideias sem encaixamentos sintáticos ; redução de estruturas subordinadas, ao lado do aumento na frequência de estruturas coordenadas e absolutas ; redução no uso de conectores para expressar relações lógicas essenciais à construção do texto ; emprego equivocado de operadores que não estabelecem relações lógicas coerentes entre ideias do texto ; fragmentação frasal que compromete a estrutura lógico-gramatical.
Uma estrutura que vem se ampliando na escrita dos participantes, refletindo hábitos da oralidade, é a duplicação de um dos elementos da frase, ou seja, a retomada anafórica de um elemento topicalizado, como por exemplo :"Mesmo com a época da censura, o ser humano ELE pode se expressar de várias formas..."; “ A cultura ELA vem trazida de várias épocas... ” ; “ Arte e cultura, ELAS \textbf{estão} relacionadas em uma só ” ; “ A cultura ELA vem enraizada a pessoa... ” ; “ As ações educativas ELAS podem ser consideradas boas ou ruins. ” Também em relação ao emprego dos sinais de pontuação, observaram -se distintos tipos de desempenho : participantes que dominam a funcionalidade da pontuação para a construção de textos fluentes e coesos, por um lado, e um grande número de participantes que apresentam precariedade em sua utilização, chegando ao ponto de elaborar parágrafos sem qualquer marca interna de pontuação para separar os períodos.
São os seguintes os tipos de desvios encontrados : - emprego de vírgula para separar o sujeito e o predicado : “ Os teatros, festas de rua, carnaval, elevam nosso conhecimento ” ( no lugar de “ Os teatros, festas de rua, carnaval elevam nosso conhecimento ” ) ;"que governos e municípios, façam a inclusão ” ( no lugar de “ que governos e municipios façam a inclusão ” ) ; - ocorrência de apenas uma das vírgulas para separar uma palavra, uma expressão ou uma oração encaixada : “ A liberdade artística que \textbf{está} presente em nossa Constituição 88, nos calça ” ( no lugar de “ A liberdade artística, que \textbf{está} presente em nossa Constituição 88, nos calça ) ” ; - emprego de ponto no lugar de vírgula para separar orações de um mesmo período : “ criando debate em espaços públicos.
Além de convidar os críticos a mostrarem o que é arte! ” ( no lugar de “ criando debate em espaços públicos, além de convidar os críticos a mostrarem o que é arte! ” ) ; - ausência de vírgula ou conjunção para separar elementos de uma enumeração :"As cidades inteligentes são caracterizadas por inteligência artificial, desenvolvimento sustentável conservação dos espaços verdes"( no lugar de"As cidades inteligentes são caracterizadas por inteligência artificial, desenvolvimento sustentável e conservação dos espaços verdes") ; - ausência de vírgula para separar oração adjetiva explicativa ou utilização inadequada para separar oração adjetiva restritiva :"As cidades inteligentes que são lugares em que a tecnologia avançada é utilizada..."( no lugar de"As cidades inteligentes, que são lugares em que a tecnologia avançada é utilizada,...") ; - emprego indevido do sinal de dois pontos : “ percebe -se um círculo onde \textbf{estão} infinitamente ligadas : a cultura e a arte ” ( no lugar de “ percebe -se um círculo onde \textbf{estão} infinitamente ligadas a cultura e a arte ” ) ; - ausência de pronome no processo de referenciação : “ censura a arte ou até mesmo impulsiona ” ( no lugar de “ censura a arte ou até mesmo a impulsiona ” ) ; - emprego de vírgula em lugar do ponto-final : “ são muitos os agravamentos em relação à tal situação, ” III - Aspectos Morfossintáticos e Vocabulares : o domínio das convenções morfossintáticas apresentou resultados que evidenciam, claramente, alguns aspectos problemáticos no desempenho dos participantes.
Assim como na análise dos desvios textuais, pode -se observar uma distinção entre três grupos de participantes : o grupo que revela domínio das convenções morfossintáticas ; o grupo que apresenta um domínio regular, com desvios pontuais ; e o grupo que revela domínio precário dessas convenções.
Pode -se destacar a regência, tanto verbal quanto nominal, como o aspecto com maior número de desvios : a ) Regência nominal e verbal : Na maioria dos textos, destacam -se os desvios relacionados à noção de"regência", ou seja, à relação de subordinação que se estabelece entre as palavras por meio, ou não, de preposições, refletindo uma tendência na modalidade escrita da língua.
O desvio mais frequente foi a falta do sinal indicativo da crase, em contextos em que seria obrigatória, ou, ao contrário, um uso indevido em contextos em que o princípio não se aplica.
Os dois tipos de desvios revelam que o usuário não tem consciência de que, sob a forma do termo “ a ”, se pressupõe a presença de uma preposição “ a ”, exigida pela regência do termo anterior.
Exemplos : “ que se contrapõem a essência das artes ” ( no lugar de “ que se contrapõem à essência das artes. ” ) ; “ Esse ato pode levar a naturalização"( no lugar de “ Esse ato pode levar à naturalização") ; “ Além de convidar os críticos a mostrarem o que é arte à fazendo! ” ( no lugar de “ Além de convidar os críticos a mostrarem o que é arte a fazendo! ” ).
Outro desvio relacionado à regência verbal e à nominal, encontrado frequentemente nas questões, é a ausência de preposição antes de pronome relativo, refletindo um processo generalizado na modalidade oral da língua, em situações de registro informal.
Exemplo :"a cidade que todos querem morar é..."( no lugar de"a cidade em que todos querem morar é...") b ) Concordância nominal e verbal : Os desvios relativos à aplicação das regras de concordância verbal e nominal correspondem a fatos que têm sido amplamente estudados na modalidade oral da Língua Portuguesa.
Quanto à concordância de número, observou -se ausência de marca ( com sujeito anteposto ou posposto ), como por exemplo a ausência de acento circunflexo na forma plural do presente do indicativo do verbo “ ter ”, considerada como um desvio na concordância verbal e não na acentuação gráfica.
Exemplo :"os governantes tem procurado planejar"( no lugar de"os governantes têm procurado planejar") ; “ onde é mostrado os acontecimentos da atualidade ” ( no lugar de “ onde são mostrados os acontecimentos da atualidade ” ).
Outro desvio na aplicação da regra da concordância de número foi o aparecimento da marca de plural em verbos ou adjetivos relacionados a núcleos substantivos no singular, devido à presença de modificadores no plural, evidenciando um processo de hipercorreção.
Esse processo tem se ampliado nas situações de oralidade, caracterizando uma tendência de variação da Língua Portuguesa.
Exemplos :"a construção de teatros e cinemas revelam uma preocupação das cidades inteligentes"( no lugar de"a construção de teatros e cinemas revela uma preocupação das cidades inteligentes") ;"o aumento dos acidentes ambientais devem ser evitados ” ( no lugar de"o aumento dos acidentes ambientais deve ser evitado").
Quanto à concordância de gênero, vários casos foram observados, normalmente no âmbito de sintagmas nominais longos, em que o adjetivo \textbf{está} afastado do substantivo, conforme apontam as pesquisas sobre a modalidade oral.
Exemplos : “ sendo valorizado uma normatização de discurso que se contrapõe ” ( no lugar de “ sendo valorizada uma normatização de discurso que se contrapõe ” ) ; “ podendo essa agradar um público e desagradar outra ” ( no lugar de “ podendo essa agradar um público e desagradar outro ” ). c ) Registro da forma verbal “ é ” sem acento, provocando ambiguidade na leitura pela confluência com a conjunção aditiva"e".
Exemplo : “ A principal forma de realizar essa disseminação e criar nas escolas... ” ( no lugar de “ A principal forma de realizar essa disseminação é criar nas escolas... ” ). d ) Emprego do pronome relativo “ onde ” como relativo universal, típica marca da oralidade, em que esse pronome não é mais utilizado apenas em sua função locativa"lugar onde se realiza a ação".
Exemplo :"a aceitação não é mais a mesma, onde muitas vezes é a cultura que censura a arte"; “ situações onde ” ; “ cultura mais cotidiana, onde ”. e ) Falta de artigo definido antes de substantivo, comprometendo o paralelismo em enumerações : Exemplo :"censurar a arte e cultura é prejudicial"( no lugar de"censurar a arte e a cultura é prejudicial"). f ) Desvios relacionados à flexão verbal : Exemplos : emprego indevido do tempo ou do modo verbal : “ adaptado ” por “ adaptando ” ; eliminação de marca de infinitivo : “ ta ” por “ \textbf{estar} ”, “ \textbf{está} ” por “ \textbf{estar} ”, “ esta ” por “ \textbf{estar} ” ; “ \textbf{estás} ” por “ estas ” ; utilização de construção de “ \textbf{estar} ” + gerúndio como marca de oralidade : “ devemos \textbf{cada} vez mais ta melhorando ” ; segmentação indevida da desinência de número e pessoa de verbos : “ entender -mos ”, “ aprender -mos ”.
Assim como nas demais competências avaliadas, há diferenças entre os participantes em relação à seleção e ao emprego do vocabulário.
Foram encontrados textos com alto padrão de domínio vocabular, chegando até ao emprego de vocabulário técnico das áreas do saber abordadas nas questões : arte, cultura, meio ambiente e tecnologia.
Por outro lado, na maioria dos textos foram observados vários tipos de inadequação na seleção e na utilização do vocabulário, sendo que o principal deles foi a excessiva repetição de \textbf{certas} palavras, revelando limitação de repertório vocabular.
Além disso, observou -se a ocorrência de expressões da oralidade, revelando falta de domínio de vocabulário mais abstrato e de maior complexidade, essencial ao desenvolvimento do texto de base dissertativa.
Exemplos : seleção vocabular indevida : “ crescemos tão intrínsecos a esses usos... ” ; emprego inadequado do verbo “ ter ” no lugar de “ haver ” : “ Durante a ditadura teve muita repressão ” ; imprecisão vocabular : “ ela anda lado a lado com a existência do que conhecemos de cultura ”, “ A liberdade artística que \textbf{está} presente em nossa Constituição 88, nos calça e... ”.
A imprecisão vocabular, como o uso da palavra “ coisa ” em vários textos, demonstra o domínio precário da norma-padrão da Língua Portuguesa, com presença constante de traços de oralidade. 6.3.1.8 CONSIDERAÇÕES FINAIS SOBRE A CORREÇÃO DE LÍNGUA PORTUGUESA DAS RESPOSTAS ÀS QUESTÕES DISCURSIVAS DO COMPONENTE DE FORMAÇÃO GERAL O conjunto de textos analisados no Enade 2021 permite traçar um perfil bem detalhado do desempenho linguístico do estudante de graduação brasileiro.
Conforme ficou comprovado nos comentários dos exemplos de respostas apresentados neste relatório, uma parcela significativa dos participantes revela lacunas importantes em sua formação, principalmente no que diz respeito à estruturação e à organização de textos.
Como esta é uma competência essencial à produção escrita, que deve ser construída ao longo dos anos de escolaridade ( do ensino fundamental ao ensino superior ), sua recuperação é trabalhosa e depende de um esforço intenso de exposição a textos e de reflexão sobre os requisitos que garantem a clareza e o encadeamento das ideias.
A competência relacionada ao domínio do vocabulário, que deve ser diversificado e apropriado ao padrão formal da língua, também apresentou um desempenho insatisfatório pela maioria dos participantes, que se limitam ao senso comum, com vocabulário genérico e vago, o que é inadequado para atender ao perfil de um graduando.
As outras duas competências avaliadas ( convenções ortográficas, convenções morfossintáticas ) apresentam desvios pontuais e são de mais fácil recuperação, com exceção dos desvios morfossintáticos que já \textbf{estão} incorporados na modalidade oral informal da Língua Portuguesa, configurando processos de variação linguística e, talvez, antecipando futuras mudanças no padrão da modalidade escrita.
Com base na avaliação do desempenho linguístico aqui relatada, conclui -se que é preciso aperfeiçoar a formação superior, dando aos estudantes mais acesso às ferramentas necessárias para que sejam bem-sucedidos nas diferentes situações de sua vida cidadã – profissional, cultural e política.
Essa é uma contribuição efetiva que os resultados do Enade podem oferecer às instituições de ensino superior do país.
Para finalizar, vale observar que as duas questões discursivas atenderam às exigências de avaliação, tanto do domínio discursivo dos estudantes, quanto do seu conhecimento sobre temas relacionados ao conhecimento acadêmico, o que pode ser extremamente útil para que os cursos de ensino superior possam contribuir na formação de seus graduandos.
Se o objetivo é avaliar a competência discursiva dos formandos brasileiros, a elaboração de uma resposta de base dissertativo-argumentativa ( Questão 1 ) e de uma resposta de base expositiva ( Questão 2 ) certamente ofereceu material relevante para essa análise, por incentivar a reflexão e a organização do pensamento por meio de uma sequência textual coerente e complexa. 6.3.2 COMPONENTE DE CONHECIMENTO ESPECÍFICO Na parte da prova relativa às Questões Discursivas do Componente de Conhecimento Específico ( Tabela 6.17 ), observa -se que a Média foi mais alta do que para as questões discursivas do Componente de Formação Geral.
Enquanto, no Componente de Formação Geral, a Média para estudantes de Música ( Licenciatura ) de todo o Brasil foi 27,6, na parte de Conhecimento Específico, a Média foi 48,1.
A maior Média deste componente foi obtida pelos estudantes da região Sul ( 53,8 ), e a menor, pelos da região Norte ( 39,9 ).
As demais regiões obtiveram Média igual a 40,6 ( Nordeste ), 51,3 ( Sudeste ) e 48,2 ( Centro-Oeste ).
Quanto à variabilidade das notas, o Desvio padrão de todo o Brasil foi 24,7.
O maior Desvio padrão foi encontrado na região Centro-Oeste ( 25,1 ), e o menor, na região Sul ( 21,9 ).
A maior nota Máxima, 100,0, foi obtida em três regiões.
Nas demais regiões, as notas Máximas foram : 88,3 ( Norte, a mais baixa ) e 93,3 ( Centro-Oeste ).
A nota Mínima zero foi obtida por, pelo menos, um estudante em todas as regiões do Brasil.
A Mediana do Brasil foi 51,7.
Nas regiões Norte e Nordeste, a Mediana foi 43,3 ( as menores ), na região Centro-Oeste foi 50,0, e nas regiões Sudeste e Sul, 56,7 ( as maiores ).
No Gráfico 6.10, representa -se a distribuição das notas nas Questões Discursivas no Componente de Conhecimento Específico.
A moda dessa distribuição ocorre no intervalo ( 50 ; 60 ], com 16,3 \% do total de participantes.
A análise de \textbf{cada} uma das Questões Discursivas de Conhecimentos Específicos será feita a seguir. 6.3.2.1 ANÁLISE DA QUESTÃO DISCURSIVA 3 DO COMPONENTE DE CONHECIMENTO ESPECÍFICO Na Questão 3, cujos resultados aferidos se encontram descritos na Tabela 6.18, a Média dos estudantes de todo o Brasil foi 46,3.
O desempenho dos estudantes nessa questão foi inferior ao da Questão 4 e superior ao da Questão 5.
A menor Média nessa questão foi obtida pelos estudantes da região Nordeste ( 38,8 ), enquanto a maior Média foi obtida na região Sul ( 51,3 ).
Quanto à variabilidade das notas, o Desvio padrão de todo o Brasil foi 29,5.
O maior Desvio padrão foi obtido na região Centro-Oeste ( 30,6 ), enquanto o menor foi obtido na região Sul ( 27,1 ).
A nota Máxima, 100,0 pontos, foi alcançada por, pelo menos, um estudante em \textbf{cada} região.
A Mediana do Brasil foi 50,0, para maioria das regiões ; nas regiões Norte e Nordeste foi 30,0.
A nota Mínima zero foi a mesma em todas as regiões do Brasil.
No Gráfico 6.11, mostra -se a distribuição das notas na Questão Discursiva 3, do Componente de Conhecimento Específico da Área de Música ( Licenciatura ).
Observa -se que 10,7 \% dos estudantes deixaram a resposta à questão em branco e que 4,2 \% receberam nota zero.
A moda dessa distribuição ocorre no intervalo ( 40 ; 50 ] que, seguido do intervalo ( 70 ; 80 ], que apresentam as maiores frequências, 27,6 \% e 20,6 \% do total de participantes, respectivamente. 6.3.2.2 COMENTÁRIOS SOBRE AS RESPOSTAS À QUESTÃO DISCURSIVA 3 A Questão Discursiva 3 atendia às Diretrizes Curriculares da prova de componentes específicos da Área de Música, Modalidade Licenciatura, no âmbito do Exame Nacional de Desempenho dos Estudantes ( Enade ), edição 2021, conforme a portaria no 397, de 23 de agosto de 2021.
Em particular, a questão possibilitou observar o comprometimento dos estudantes com a promoção do acesso à educação inclusiva.
Nesse sentido, o enunciado dessa questão discursiva, que era claro e objetivo, incentivou o estudante a buscar formas de inclusão efetiva de pessoas com deficiências, nesse caso a deficiência visual.
Tal questão se apresentou como transversal aos conteúdos contemplados pela Área da Música, tomando como referência o Art. 7o das Diretrizes da Área, uma vez que o processo de inclusão precisa ser levado em conta na aprendizagem musical.
O espaço para resposta de 15 linhas mostrou -se pequeno, já que a questão se subdividia em dois itens, especialmente porque o segundo previa a descrição de uma atividade de ensino que incluísse a apresentação de material didático e da avaliação.
A questão pode ser considerada de dificuldade média, por requerer aspectos reflexivos, descritivos e a apresentação de material didático e de instrumentos de avaliação.
O padrão de resposta abarcava respostas bem diversas, apesar de centradas na deficiência visual, mas tal diversidade não foi encontrada no conjunto de respostas dos estudantes.
No item ‘ a ’, esperava -se que o respondente apontasse dificuldades de aprendizagem na Área de Música para estudantes com deficiência visual e que discutisse a forma como tais dificuldades podem interferir no processo de aprendizagem.
Nesse item, os estudantes limitaram -se, na maioria das vezes, a indicar uma dificuldade de aprendizagem.
No item ‘ b ’, o estudante deveria descrever uma atividade de ensino de Música, sendo esperada pelo padrão de respostas a apresentação de estratégias, mostrando adaptações necessárias para atender ao deficiente visual.
As respostas dadas pelos estudantes mostraram -se bastante sucintas e, em muitos casos, eles desconsideraram a solicitação de apresentar os materiais didáticos a serem empregados e a forma de avaliação.
A linguagem adotada pelos alunos foi, em geral, adequada.
Essas respostas, em geral, corresponderam ao senso comum de que as dificuldades encontradas por estudantes de Música com deficiência visual centram -se na impossibilidade da leitura.
Nesse sentido, a solução mais apresentada pelos respondentes foi indicar a necessidade de material de estudo em braile.
Apesar da pertinência, nas respostas, os estudantes não mostraram esmero em discutir de forma mais reflexiva a questão, nem demonstraram uma compreensão mais ampla do processo de inclusão.
Assim, apesar de não \textbf{estarem} incorretas, as respostas foram superficiais, talvez indicando insuficiência a respeito das discussões sobre inclusão durante o processo formativo nos cursos de Licenciatura.
De De modo geral, as respostas \textbf{estavam} incompletas e/ou atendendo a apenas um dos aspectos solicitados pela questão.
A Questão 3 teve 10,7 \% de respostas em branco.
As consideradas fracas ( notas de 00 a 30 ) foram 28,2 \% ; 31,8 \% foram medianas ( notas de 35 a 70 ) ; e 29,2 \%, boas ( notas de 75 a 100 ).
De forma geral, as fracas foram aquelas em que os estudantes se limitaram a identificar uma dificuldade de aprendizagem em Música que um com algum grau de deficiência visual pode apresentar, entretanto, em alguns casos, sem comentar como a dificuldade interfere no processo de aprendizagem ( resposta relativa ao item ‘ a ’ da questão ).
As medianas foram aquelas em que os estudantes responderam ao item ‘ a ’, mas deixaram o item ‘ b ’ incompleto, aquele que previa a descrição de uma atividade de Música e a indicação de material didático e da avaliação.
Nas respostas consideradas boas, observou -se um maior cuidado no \textbf{detalhamento} ao que foi solicitado nos itens ‘ a ’ e ‘ b ’.
A correção dessa questão apontou para a necessidade de uma atenção especial para o tema da inclusão, abrangendo estratégias para \textbf{superação} das necessidades de pessoas especiais, durante a formação do professor de Música.
Cabe registrar que o fraco desempenho ocorreu mesmo com uma correção bastante flexível, já que as respostas foram, no geral, sucintas e pouco detalhadas, até mesmo por conta do pouco espaço para que fossem mais descritivas e explicativas.
Foi identificada relação da Questão Discursiva 3 com as Questões objetivas de números 22 e 33 que também tratavam sobre educação inclusiva, o que reforçou a importância da temática na formação dos professores. 6.3.2.3 ANÁLISE DA QUESTÃO DISCURSIVA 4 DO COMPONENTE DE CONHECIMENTO ESPECÍFICO Na Tabela 6.19, constam as informações relativas à Questão 4 do conjunto de Questões do Componente de Conhecimento Específico.
O desempenho dos estudantes nessa questão foi superior ao das Questões 3 e 5.
A Média geral do Brasil foi 60,0, sendo a maior Média registrada na região Sul ( 65,8 ), e a menor, na região Nordeste ( 51,7 ).
A nota Máxima ( 100,0 ) foi atingida por, pelo menos, um concluinte em todas as regiões.
A Mediana em todo o Brasil foi 70,0, o mesmo valor foi obtido na região Centro-Oeste.
As regiões Norte e Nordeste obtiveram valor inferior ao do Brasil, 60,0.
E as regiões Sudeste e Sul obtiveram valor superior ao Brasil, 80,0.
Foi zero a nota Mínima de todas as regiões.
No Gráfico 6.12, \textbf{está} representada a distribuição de notas da Questão Discursiva 4, do Componente de Conhecimento Específico.
Essa distribuição tem moda no intervalo ( 90 ; 100 ] com 22,5 \% do total.
Nessa questão, 2,2 \% obtiveram nota zero e 19,2 \% deixaram a resposta em branco.
O intervalo ( 70 ; 80 ] caracteriza -se como um máximo local, com cerca de 13,2 \% do total. 6.3.2.4 COMENTÁRIOS SOBRE AS RESPOSTAS À QUESTÃO DISCURSIVA 4 A Questão Discursiva 4 relacionava -se, particularmente, ao item VII, artigo 5o, da Portaria no 397, de 23 de agosto de 2021, que dispõe sobre Diretrizes da Área de Música, Modalidade Licenciatura, no âmbito do Exame Nacional de Desempenho dos Estudantes ( Enade ), edição 2021.
A questão era adequada, quanto à referência ao perfil esperado do concluinte pelas Diretrizes, que aponta a característica de ser “ competente, expressivo e criativo no fazer musical e em sua prática docente ”.
A questão tratava das metodologias dos educadores musicais da primeira geração dos métodos ativos, que tiveram e têm, até os dias de hoje, grande influência sobre os processos de ensino-aprendizagem musical.
Nesse sentido, a questão se articulava com os diversos conteúdos previstos pela referida Portaria em seu artigo 7o.
O enunciado \textbf{estava} claro e objetivo.
Apesar de delimitar a temática de ensino e aprendizagem conforme preceitos dos educadores da primeira geração dos métodos ativos, dava liberdade de escolha entre eles para servir como inspiração na elaboração de uma atividade musical direcionada para estudantes de 8 a 10 anos e que tratasse de algum parâmetro sonoro especificado nesse enunciado ( altura, duração, intensidade e timbre ).
Apesar de \textbf{estar} subdividida em quatro itens, o espaço para a resposta foi considerado suficiente.
A questão poderia ser considerada fácil, se comparada às demais, sob o ponto de vista de que saber elaborar uma atividade musical é a base da formação de um licenciado em Música.
O padrão de resposta previa a indicação de um educador musical da primeira geração dos métodos ativos, como fonte de inspiração para a elaboração de uma atividade a ser descrita na resposta.
Porém, com frequência, os estudantes indicaram um educador da segunda geração.
Esse fato não desqualificava a resposta como um todo, pois os outros itens foram avaliados observando a coerência entre objetivos, estratégias didáticas e avaliação, mesmo não correspondendo ao educador musical indicado ou correspondendo a um educador que não fosse da primeira geração.
A linguagem adotada pelos alunos foi adequada do ponto de vista dos termos específicos da Área.
Chamou a atenção a referência constante ao educador musical Émile Jaques-Dalcroze, pertencente à primeira geração ( vale notar que a grafia do nome do educador, com muita frequência, apareceu de forma equivocada ), e ao educador musical da segunda geração dos métodos ativos Raymond Murray Schafer.
Estas foram as escolhas de abordagem usadas com mais frequência para responder à questão.
Os planos de aula apresentaram diversidade, e o principal equívoco encontrado foi a confusão entre educadores da primeira e da segunda gerações dos métodos ativos, ou a indicação de outros educadores musicais não pertencentes a essas duas gerações.
A Questão Discursiva 4 teve como resultado 6,6 \% de respostas consideradas fracas ( notas de 00 a 30 ) ; 27,4 \%, medianas ( notas de 35 a 70 ), e 35,7 \%, boas ( notas de 75 a 100 ), evidenciando um desempenho melhor do que o da Questão 3.
No entanto, o percentual de respostas em branco foi maior, 19,2 \%.
As respostas medianas, em geral, \textbf{estavam} incompletas em relação ao esperado pelo padrão de resposta, nas quais os estudantes apresentavam pouco \textbf{detalhamento} na descrição de objetivos, estratégias didáticas e/ou mecanismos de avaliação de aprendizagem.
Em muitos casos, ainda, se equivocaram na resposta ao item ‘ d ’, quando deveriam indicar um educador musical da primeira geração dos métodos ativos.
Nas respostas consideradas boas, cumpriram o esperado, com um \textbf{detalhamento} maior do que as outras e apresentando coerência entre os itens da questão.
Embora em muitas respostas demonstrassem equívocos sobre os educadores da primeira geração dos métodos ativos, os estudantes foram capazes de apresentar um plano de atividades com coerência interna, ou seja, com uma boa articulação entre objetivos, estratégias didáticas e mecanismos de avaliação.
Os educadores dos métodos ativos para a educação musical, das chamadas primeira e segunda gerações, são de fundamental importância para as metodologias educativas adotadas na contemporaneidade.
Essa foi a única questão em toda a prova, dentre as discursivas e objetivas, que avaliou esse tipo de conhecimento. 6.3.2.5 ANÁLISE DA QUESTÃO DISCURSIVA 5 DO COMPONENTE DE CONHECIMENTO ESPECÍFICO Na Tabela 6.20, constam as informações relativas à Questão 5 do conjunto do Componente de Conhecimento Específico.
O desempenho dos estudantes de todo o Brasil nesta questão foi inferior ao da Questão Discursiva 3 e ao da Discursiva 4, desse componente.
A nota Média dos estudantes de todo o Brasil foi 38,1.
A maior Média foi registrada na região Sul ( 44,2 ), enquanto a menor Média foi registrada na região Norte ( 28,1 ).
Quanto à variabilidade das notas, o Desvio padrão dos estudantes do Brasil, como um todo, foi 28,0.
Enquanto o maior desvio foi encontrado na região Centro-Oeste ( 30,0 ), o menor foi encontrado na região Sul ( 25,9 ).
A Mediana para o Brasil e em todas as regiões foi 40,0.
Para o conjunto de estudantes de Música ( Licenciatura ) do Brasil, a nota Máxima foi 100,0, obtida por, pelo menos, um estudante em todas as regiões.
A nota Mínima do Brasil foi zero, o mesmo valor encontrado em todas as regiões.
No Gráfico 6.13, é representada a distribuição das notas da Questão Discursiva 5 do Componente de Conhecimento Específico.
Essa distribuição tem moda no intervalo ( 30 ; 40 ], que corresponde a 32,1 \% dos respondentes.
Os que obtiveram nota zero correspondem a 24,4 \%, a soma daqueles que deixaram a resposta à questão em branco com os que tendo respondido não receberam pontuação. 6.3.2.6 COMENTÁRIOS SOBRE AS RESPOSTAS À QUESTÃO DISCURSIVA 5 A Questão Discursiva 5 atendeu às Diretrizes de prova e de componentes específicos da Área de Música, \textbf{estando} de acordo com os seguintes itens listados pelo Art. 7o das No Gráfico 6.13, é representada a distribuição das notas da Questão Discursiva 5 do Componente de Conhecimento Específico.
Essa distribuição tem moda no intervalo ( 30 ; 40 ], que corresponde a 32,1 \% dos respondentes.
Os que obtiveram nota zero correspondem a 24,4 \%, a soma daqueles que deixaram a resposta à questão em branco com os que tendo respondido não receberam pontuação.
Diretrizes da Área : I - Elementos estéticos, históricos e filosóficos da música ; III - Elementos perceptivos, teóricos, analíticos e estruturais da música ; e IV - Prática musical : elementos de interpretação, repertórios, estilos e gêneros.
A questão apresentava clareza em seu enunciado, porém, as ilustrações presentes ( partituras ) podem ter influenciado a tomada de decisões dos estudantes na resposta ao item ‘ a ’, em que deveriam escolher qual das partituras seria a mais adequada à questão apresentada.
Nesse sentido, encontramos interpretações diferentes da esperada pelo padrão adotado.
A questão pode ser considerada difícil por esse motivo, além de não ser, necessariamente, um conhecimento comum ao perfil dos licenciados que têm origens, experiências e práticas musicais muito diversas.
O enunciado apresentava duas partituras de um mesmo trecho musical com edições diferentes.
As imagens mostravam a data em que as edições foram realizadas, a primeira em 1950 e a segunda em 2019.
A tarefa do estudante seria identificar qual delas se aproximaria melhor da forma de notação original do compositor da obra J.
S.
Bach.
A grande maioria dos respondentes indicou a primeira edição, provavelmente por esta ser a mais antiga conforme a data que constava da imagem, porém, a resposta correta seria a indicação da segunda partitura.
A datas das edições, possivelmente, induziram os estudantes ao erro e, talvez por isso, tenha sido grande o número de respostas em branco ou equivocadas, em relação ao item ‘ a ’.
Não houve grande diversidade nas respostas, e o principal erro cometido foi na indicação da partitura.
Mesmo assim, foram consideradas na correção as observações dos estudantes que indicaram, corretamente, qual seria a diferença entre as duas versões da música.
Os estudantes não tiveram bom desempenho nesta questão : 13,3 \% das respostas foram consideradas fracas ( notas de 00 a 30 ), 51,8 \% das respostas receberam notas medianas, ( notas de 35 a 70 ) e apenas 13,5 \% foram respostas boas ( notas de 75 a 100 ).
Foi a questão com a maior percentagem de respostas em branco, 21,3 \%.
Nas fracas, em geral, se equivocaram na escolha da partitura, mas conseguiram indicar alguma diferença entre elas.
Nas medianas, em muitos casos, também se equivocaram em relação à escolha da partitura, mas trouxeram maior \textbf{detalhamento} sobre as diferenças presentes nas duas edições e sobre a forma de interpretação da obra.
Nas respostas consideradas boas, indicaram corretamente a partitura, detalharam as diferenças entre as edições e a forma como essas diferenças interferem na interpretação da obra.
Por meio das respostas, constatou -se uma deficiência na compreensão das características da música barroca europeia.
Os estudantes buscaram responder à questão pautando -se na data da edição da partitura e não nas características musicais indicadas por elas.
Apesar disso, souberam distinguir as particularidades da grafia e dos sinais de dinâmica apresentados.
A correção buscou considerar os conhecimentos apresentados pelos estudantes no \textbf{detalhamento} das respostas.
Mesmo quando a indicação da partitura foi equivocada, privilegiou -se a argumentação sobre as diferenças da grafia das edições e a implicação dessas diferenças na interpretação da obra.
A prova do Enade/2021 apresentou muitas questões objetivas, além desta discursiva, que se atinham a elementos técnicos, como a identificação de tonalidades, fraseologias, etc.
A prova, no que se refere à parte musical, \textbf{está} mais técnica do que reflexiva.
Talvez pelo caráter mais prático-reflexivo do que técnico dos cursos de Licenciatura, tenham sido observadas muitas respostas em branco e muitas notas baixas nessa questão. 6.3.2.7 CONSIDERAÇÕES FINAIS SOBRE AS QUESTÕES DISCURSIVAS DO COMPONENTE DE CONHECIMENTO ESPECÍFICO O resultado da correção das Questões Discursivas do Componente de Conhecimento Específico revelou aspectos importantes que merecem ser apreciados pelos cursos de Licenciatura em Música.
A Questão Discursiva 3 trouxe o debate sobre a importância da inclusão de pessoas com deficiência visual nos processos de ensino-aprendizagem musical.
Os estudantes apresentaram respostas sucintas e calcadas no senso comum, mostrando a necessidade de um aprofundamento em tal temática na formação de professores de Música, no sentido de buscar o comprometimento na promoção do acesso à educação inclusiva de qualidade.
A Questão de número 4 foi a de melhor aproveitamento.
A despeito de \textbf{certa} confusão na identificação dos educadores musicais da primeira geração dos métodos ativos, o desempenho foi satisfatório no que diz respeito à capacidade para formular propostas de ensino musical.
Tal resultado demonstrou capacidade dos estudantes em integrar teoria e prática nos processos de ensino e aprendizagem musicais.
No entanto, o desenvolvimento e a aplicação de estratégias de avaliação na prática docente em Música mostraram -se frágeis.
Já a Questão de número 5 demonstrou a necessidade de um maior conhecimento sobre as características da música barroca europeia.
Vale destacar que as duas primeiras Questões do Componente de Conhecimento Específico ( 3 e 4 ) são de extrema importância para a formação do Licenciando em Música, considerando -se a diversidade de espaços onde poderá atuar.
A última foi a de pior desempenho, talvez por ser mais específica e privilegiar determinada cultura, dada a variedade dos saberes estético-musicais dos respondentes.
Como recomendação final, a partir do observado por meio da correção das respostas dos Licenciandos em Música, destaca -se a necessidade de ampliação e/ou aprofundamento nos currículos de Música a respeito de questões relacionadas à atenção e à empatia às diferenças, em particular no que se refere à inclusão de pessoas com deficiências.

% matched lemmas: cada, certo, detalhamento, estar, superação, torno
\end{textsample}
